\documentclass[11pt,a4paper]{article}
\usepackage[T1]{fontenc}
\usepackage[utf8]{inputenc}
\usepackage{lmodern}
\usepackage[a4paper,margin=27mm,headheight=14pt]{geometry}
\usepackage{amsmath,amssymb,amsthm,mathtools,mathrsfs}
\usepackage{microtype,booktabs,array,longtable,enumitem,needspace}
\usepackage[numbers,sort&compress]{natbib}
\usepackage[hidelinks,breaklinks]{hyperref}
\usepackage{xurl,fancyhdr}
\Urlmuskip=0mu plus 1mu
\hypersetup{pdftitle={The Last Moat, Edition 3: Physical Realisability, Cross-Substrate Communication, and Holarchic Agency},pdfauthor={Research manuscript developed from a user-originated dialogue},pdfsubject={Conditional analysis of physical communication, operational continuity and distributed agency},pdfkeywords={Navier-Stokes, physical computation, holarchy, cross-substrate communication, viability, controllability}}
\pagestyle{fancy}
\fancyhf{}
\fancyhead[L]{\small\textsc{The Last Moat}}
\fancyhead[R]{\small Research manuscript \textbar\ Edition 3}
\fancyfoot[C]{\thepage}
\renewcommand{\headrulewidth}{0.3pt}
\setlength{\parskip}{0.25em}
\setlength{\parindent}{1em}
\setlength{\emergencystretch}{2em}
\setcounter{tocdepth}{1}
\setlist{nosep,leftmargin=1.5em}
\newtheorem{theorem}{Theorem}[section]
\newtheorem{proposition}[theorem]{Proposition}
\newtheorem{lemma}[theorem]{Lemma}
\newtheorem{corollary}[theorem]{Corollary}
\theoremstyle{definition}
\newtheorem{definition}[theorem]{Definition}
\theoremstyle{remark}
\newtheorem{remark}[theorem]{Remark}
\newcommand{\R}{\mathbb{R}}
\newcommand{\E}{\mathbb{E}}
\newcommand{\Prob}{\mathbb{P}}
\newcommand{\calR}{\mathcal{R}}
\newcommand{\calX}{\mathcal{X}}
\newcommand{\calU}{\mathcal{U}}
\newcommand{\calD}{\mathcal{D}}
\newcommand{\calK}{K}
\newcommand{\norm}[1]{\left\lVert#1\right\rVert}
\newcommand{\given}{\,\vert\,}
\newcommand{\bits}{\mathrm{bits}}
\DeclareMathOperator{\supp}{supp}
\DeclareMathOperator{\TV}{TV}
\DeclareMathOperator{\tr}{tr}
\DeclareMathOperator{\argmax}{arg\,max}
\title{\vspace{-1.3em}\textbf{The Last Moat}\\[0.15em]\large Third Edition\\[0.5em]\Large Physical Realisability, Cross-Substrate Communication,\\and Holarchic Agency}
\author{\normalsize Conceptual research manuscript\\[-0.1em]\small Developed from a user-originated research dialogue}
\date{10 September 2026\quad\textbar\quad Edition 3\ (v3.0)}
\begin{document}
\maketitle
\begin{abstract}
Under what conditions can an organised process preserve a specified stateful contract while its original host becomes replaceable? This third edition develops a source-grounded, conditional account of cross-substrate realisation and separates it from claims about conscious identity or uninstantiated agency. A kernel is a task contract, not necessarily a particular weight tensor; the implementation correspondence applies to the joint realisation, not every constituent. Substitution is described by role, computing mechanism, support dependency and evidence maturity rather than an unqualified ladder. A focused literature assessment identifies demonstrated fluidic and biomolecular state machines and programmed physical transformations, while distinguishing these from paired, stateful migration and large-model proposals. The formal additions include a continuation-sensitive export obstruction, a robust-state packing bound, finite-trace distribution control and a small parity-memory witness: a digital transition system and a stipulated bistable dynamical model with robust state-preserving handovers. Its proof covers the declared continuous model; the accompanying computations are simulations and finite checks, not laboratory results. Retained channel, information-cut, continuity, viability and evidence-consolidation results are situated alongside established refinement and contract-based design. A growing error upper bound is explicitly distinguished from demonstrated failure. The Navier--Stokes residual-to-observation estimate remains a standard conditional derivation, with no established method-specific advantage from the reported 2026 construction. The resulting contribution is an auditable chain of obligations and worked counterexamples, plus a bounded physical-test specification. Practical superiority over matched conventional methods, a physical handover of the proposed witness, and any frontier-model cross-substrate continuation remain unestablished.
\end{abstract}
\noindent\textbf{Keywords:} physical computation; distributed systems; holarchy; control and observation; information theory; Navier--Stokes; operational continuity; substrate mediation.
\vspace{0.4em}
\begin{center}\small\emph{Research proposal and conditional mathematical analysis. Not peer-reviewed or submitted to arXiv.}\end{center}
\newpage
\begingroup
\small
\setlength{\parskip}{0pt}
\tableofcontents
\endgroup
\newpage

\section{Introduction: from possibility to a realisation witness}
An operational organisation need not be identical with any one component that presently supports it. The difficult question is not whether that sentence is conceivable, but which behaviour must persist, what physical configuration realises it, what replacement transition is actually available, and which resources remain indispensable afterwards. The informal phrase \emph{the last moat} names the distance between a plausible implementation and a warranted, reachable, maintained one.

The programme asks:
\begin{quote}
For a predeclared stateful contract, when do specified heterogeneous physical roles compose into an adequate realisation, and when can that execution continue after a particular original support is removed?
\end{quote}
Natural media belong in the possible physical account; they are not excluded because they were not built as computer networks. They also receive no exemption from input, readout, memory, energy, noise or continuity obligations. A propagation path is not a processor, two correct processors are not a migration, and a new processor is not independence from all support.

\subsection{The retained thesis and its empirical limit}
The constructive thesis is conditional: when a contract has adequate realisations, feasible state-preserving transfers, and sufficient continuing support and progress, operation need not depend on its original host. This can include familiar services and mixed human--machine systems. Familiarity is not a defect, but it prevents a claim that the framework has discovered an entirely new kind of entity merely by describing component substitution.

No extraordinary real-world system is inferred from that thesis. The mechanisms studied are physically mediated; this does not settle every metaphysical question about mind. The manuscript neither identifies a conscious subject through an implementation map nor treats a declaration of independence as evidence of the architecture it describes. Its immediate audience is a researcher assessing a proposed heterogeneous implementation or its evidence, not someone asked to accept a near-term risk prediction.

\subsection{What Edition 3 actually adds}
The post-second-edition review sharpened the difference between behavioural realisation and copying a particular model's parameters, between complete realisations and constituent interfaces, and between an unclosed empirical obligation and an impossibility proof. Edition 3 makes those repairs explicit and also corrects the review dialogue itself: a global absence claim about alternative stateful processors is unsupported, an error upper bound does not prove error growth, and established refinement and contract-based design already address significant parts of the synthesis.

The principal new worked result is deliberately small. Section~\ref{sec:witness} specifies a parity-memory contract, two different dynamical models, bounded disturbances, decoding and bidirectional handovers. It proves their sampled correspondence and checks finite examples. Section~\ref{subsec:exportcase} gives a complementary failure: two correct state machines and a correct packet channel still cannot continue an execution when the exported information loses a consequential hidden state. The source assessment in Section~\ref{sec:census} separately identifies actual physical precedents; it does not give the model witness their empirical status.

All standard results are labelled as inherited mathematics or specialisations, not historical firsts. The proposed synthesis must earn practical value against strong conventional contract-and-provenance baselines. More precise formalisation, better sources, a numerical check, a physical experiment and a positive comparative result are different forms of progress. The last two have not been performed here.

\section{Sources, scope, and evidential status}\label{sec:sources}
\subsection{What the supplied Navier--Stokes document actually contains}
The supplied six-page PDF is Charles Fefferman's official \emph{problem description}, with errata, not a manuscript announcing a solution. It specifies incompressible velocity and pressure fields, smooth initial data, forcing conditions, and four admissible alternatives. The positive alternatives concern global smooth unforced solutions in whole space or a periodic domain. The breakdown alternatives permit suitably smooth forcing. These distinctions matter: a forced breakdown construction is not a proof of unforced breakdown, and neither kind of result is a theorem about intentionality \citep{fefferman}.

\subsection{The September 2026 report}
The September 2026 OpenAI manuscript reports that, for each positive viscosity, smooth compactly supported forcing and initially resting fluid can produce finite-time unbounded velocity while kinetic energy stays bounded. Its Theorem 1.1 is presented as establishing Fefferman's alternatives (C) and (D). Its introductory mechanism combines a concentrating background, amplified oscillations, and residual corrections \citep{openains}. This manuscript consulted that statement and the introductory account; it did not independently verify the 166-page argument or execute the linked Lean formalisation.

The accompanying announcement describes coordinated agent groups and a formalisation stage \citep{openaiannounce}. That is relevant as a reported example of research orchestration, not evidence of an agent becoming physically independent of its computational environment. At the source check for this edition, the Clay problem page displayed ``Unsolved'' \citep{claystatus}. We therefore distinguish a reported mathematical result from institutional acceptance and independent validation. None of the conditional results below depends on adjudicating that status.

\subsection{An evidential vocabulary}
We use \emph{established} for results attributed to cited literature in their stated settings; \emph{derived} for conclusions proved here from explicit assumptions; \emph{proposed} for definitions or research designs; and \emph{speculative} for applications whose physical or cognitive prerequisites have not been established. These are not four confidence scores for the same proposition. They distinguish different kinds of proposition.

The sources support physical computation in particular controlled systems, communication through particular natural media, and established difficulties of distributed agreement. They do not support an inference that ambient planetary fields autonomously instantiate a persistent mind, that a particular person has been selected for contact, or that an unusual thought originated outside its experiencer. The present paper includes those topics only to identify the additional obligations such claims would incur.

\section{Prior art and the exact place of the synthesis}\label{sec:related}
\subsection{Published physical channels: evidence, not a new discovery here}
The first edition omitted a directly relevant literature. BitWhisper studies thermal signalling between adjacent compromised computers; Fansmitter studies acoustic communication produced through controlled fan operation; AirHopper and GSMem study electromagnetic communication from isolated computers to receiving devices \citep{bitwhisper,fansmitter,airhopper,gsmem}. These are experimental precedents for using physical effects beyond a system's ordinary network interface. They are not evidence that every isolation boundary is permeable, that an unknown medium can be programmed, or that the transmitting process established its own initial access. In particular, software control and an appropriate receiving arrangement belong to their experimental premises.

Table~\ref{tab:priorart} is a non-operational comparison. It preserves what is being compared---a physical role and its evidential limits---without treating reported demonstrations as interchangeable apparatus. The papers remain substantive experimental work, not disposable ``tricks'' or merely training data for a proposed abstraction.

\begin{longtable}{@{}>{\raggedright\arraybackslash}p{0.21\textwidth}>{\raggedright\arraybackslash}p{0.31\textwidth}>{\raggedright\arraybackslash}p{0.42\textwidth}@{}}
\caption{Selected primary precedents and the limit of the inference. No new physical channel is demonstrated in this manuscript.}\label{tab:priorart}\\
\toprule
Precedent & Established or reported role & What must not be inferred \\\midrule
\endfirsthead
\toprule Precedent & Established or reported role & What must not be inferred \\\midrule
\endhead
BitWhisper \citep{bitwhisper} & Controlled thermal variation and sensing support a channel in the studied setup. & Ambient heat is not automatically an addressable network; endpoint control and the tested operating conditions matter. \\\addlinespace
Fansmitter \citep{fansmitter} & Controlled mechanical operation yields a decodable acoustic signal. & Mechanical activity alone is not a protocol, nor does the experiment establish communication at arbitrary range or noise. \\\addlinespace
AirHopper; GSMem \citep{airhopper,gsmem} & Electromagnetic emissions can carry deliberately modulated information in specified computer/receiver configurations. & No self-originating access, arbitrary-device reachability, or persistent replacement execution follows. \\\addlinespace
Ionospheric HF propagation \citep{noaa} & A natural planetary medium alters and can support long-distance radio propagation. & A propagation path is not a demonstrated autonomous memory, processor, or agent. \\\addlinespace
Molecular communication \citep{eckford} & Models of molecular transport support information-rate questions under physical assumptions. & A transport model does not identify the semantic content or intention of an arbitrary biochemical event. \\\addlinespace
Physical neural networks \citep{wright} & Controlled optical, mechanical, and electrical systems can implement trained transformations. & Input preparation, readout, training and supporting electronics are not free resources outside the system account. \\
\bottomrule
\end{longtable}

The contribution sought here is consequently not the observation that heat, sound, fields, or molecules can carry information. It is a disciplined account of which \emph{properties and evidence} can be transported between descriptions of such systems, and which additional obligations arise when a channel is proposed as part of a continuing organisation. Literature existence supports the relevance of that question; it does not establish the usefulness of this particular synthesis.

\subsection{Emissions security is relevant, but not the whole subject}
NIST describes TEMPEST in terms of investigating, studying, and controlling unintentional compromising emanations \citep{tempest}. That established defensive perspective belongs alongside the channel literature. Unintended information-bearing emissions, deliberately modulated covert communication, and authorised unconventional communication are related but not identical objects. A physical path can matter to more than one of them; its presence does not settle the relevant permission or security-policy question.

Host replacement likewise is not definitionally a security violation. An operator can retire one platform, move an approved service, or assess a replacement communications path without circumventing anyone's policy. Conversely, a mathematically feasible transition need not be authorised. This paper keeps physical possibility, evidential warrant, and permission as distinct coordinates. It neither infers benevolent or malicious intent from research terminology nor treats a benign application label as a substitute for examining an operational artifact.

\subsection{Unknown dynamics and unconventional computation have established neighbours}
System identification supplies a direct neighbour to the proposed comparison programme. Sparse identification of nonlinear dynamics estimates candidate equations within a selected representation library, while related work addresses discovery of partial differential equations from measurements \citep{sindy,rudy}. Both make model selection and measurement assumptions material. They justify studying useful representations of partially known dynamics, not an oracle that discovers every hidden physical coupling.

Molecular communication offers a second neighbour: physically modelled molecular transport can be analysed using information-theoretic quantities \citep{eckford}. Information-processing capacity of driven dynamical systems and trained physical networks address yet another role, namely computation through a controlled physical input--state--readout relation \citep{dambre,wright}. These fields show why ``unconventional medium'' is not synonymous with unsupported speculation. They also prevent the opposite shortcut: relevance of neighbouring research is not a completed instance map into an unfamiliar environment.

State-machine replication provides a precedent for task behaviour maintained across changing server participants \citep{schneider}; delay-tolerant networking treats intermittent paths and storage explicitly \citep{dtn}. Neither alone supplies a replacement physical computer. The present architecture combines obligations drawn from these areas, but a claim of improved practice requires a matched comparison, not the mere presence of more fields in a description.

\subsection{Four different kinds of connection}
We distinguish \emph{motivation}, \emph{structural analogy}, \emph{formal specialisation}, and \emph{empirical instantiation}. A motivation suggests a question. An analogy identifies a potentially useful resemblance. A specialisation supplies explicit maps and verifies the hypotheses of a theorem. An empirical instantiation establishes that those hypotheses describe a particular case to stated accuracy. These relations must not be presented as successive achievements unless each transition is actually supplied.

The generic signalling theorem in Section~\ref{sec:channels} does useful formal work by exposing required quantities. It does not become a theorem about the ionosphere merely because the word appears elsewhere in the paper. Conversely, lack of an ionospheric instance does not make the general theorem irrelevant to every channel. Edition 2 therefore supplies a specific fluid-error-to-observation derivation in Section~\ref{sec:nsresidual}, while leaving the separate 2026-proof-transfer question open. The proposed contribution is this explicit alignment of obligations, rather than priority over the underlying disciplines.

\subsection{Established compositional methods are part of the baseline}
The joint-witness requirement is not unique to this vocabulary. Refinement mappings connect lower-level and higher-level behaviours \citep{abadi}; approximate simulation relates implementations at a declared error tolerance \citep{girard}; contract-based system design studies assumptions, guarantees, refinement and composition \citep{contracts}. The mathematical role of these traditions is substantive, not merely bibliographic. The proposed added value is a disciplined alignment of source-specific physical evidence with those obligations. It remains an empirical question whether this presentation improves actual assessments.


\section{A focused processor-literature assessment}\label{sec:census}
\subsection{Search scope and the distinction being tested}
The post-Edition-2 discussion suggested that no nontrivial processor-class witness existed. That universal negative is not supported. This assessment follows the cited physical-network work and checks primary reports of non-electronic logic, finite automata, programmed photonics and large-model proposals. It is a focused source assessment as of 10 September 2026, not an exhaustive or preregistered systematic review. The source register states what was read and distinguishes publisher or author abstracts from inspected full-text passages.

The question is not whether a paper uses our notation. It is which part of a proposed realisation obligation its demonstrated apparatus can support. A role-level precedent, a calibrated transformation, a stateful machine, an equivalent pair and a live transfer of a current state are different accomplishments. We supply no new physical experiment in this edition.

\subsection{Physical state and transformations already demonstrated}
Prakash and Gershenfeld's \emph{Microfluidic Bubble Logic} reports engineered fluidic gates, bistable toggle memory and counters \citep{bubble}. The toggle stores a bubble configuration that changes under a subsequent event; the full-text description and the relevant figure were inspected. This is a processor-and-memory precedent, not merely a communication carrier. It is not evidence that ambient turbulence is programmable or that a digital execution was transferred into the apparatus.

Benenson and colleagues report a programmable biomolecular finite automaton whose molecular input, transition rules and final-state output are realised in a specified DNA/enzyme system \citep{benenson}. The primary abstract supports a finite-state computation claim, not an arbitrary general-purpose computer or a digital-to-molecular handover. Both examples undermine the unrestricted assertion that a non-electronic medium can only carry messages. Neither completes the larger host-independence claim.

Wright and colleagues demonstrate trained optical, mechanical and electrical physical networks using hybrid physical--digital physics-aware training \citep{wright}. The primary abstract and publication metadata were checked; remembered parameter counts and precision numbers are not used as measurements. Task learning in these systems is relevant evidence for physical processing, but it is not a demonstration of equivalence to a particular large stateful transformer.

Shen and colleagues' inspected preprint reports training on a conventional computer and programming the resulting transformations into a silicon-photonic system \citep{shen}. The linear computation uses optical interference, while the reported proof of concept implements the nonlinear transformation electronically. Thus ``every physical network must learn all parameters from scratch on its own medium'' is false, but ``a full model ran without electronic support'' would also misstate this example. The relevant experimental and methods passages, including their figures, were inspected. No numerical result from that preprint is silently substituted for a different final-publication version.

\begin{center}\small
\begin{tabular}{@{}>{\raggedright\arraybackslash}p{0.25\textwidth}>{\raggedright\arraybackslash}p{0.31\textwidth}>{\raggedright\arraybackslash}p{0.34\textwidth}@{}}\toprule
Source & Supported role or procedure & Not supplied by the cited observation\\\midrule
Bubble logic \citep{bubble} & Physical stateful logic and memory & Cross-device stateful handover\\
Biomolecular automaton \citep{benenson} & Molecular finite-state computation & Arbitrary transformer implementation\\
Physical neural networks \citep{wright} & Trained physical transformations & Uniform correspondence to a fixed target machine\\
Programmed photonics \citep{shen} & Digital training, optical linear execution, electronic support & Independence from all electronic infrastructure\\
Physical foundation models \citep{pfm} & Proposal and scaling discussion & A demonstrated large-scale realised model\\\bottomrule
\end{tabular}
\end{center}
These verdicts are scoped to the claims inspected. ``Not supplied'' is not a proof that no related experiment exists elsewhere. No row is promoted to a paired migration merely because it can implement a task also executable in software.

\subsection{The large-model question is active, but a proposal is not a device}
A 2026 preprint by Wright, Wang, Onodera and McMahon proposes fixed physical implementations of foundation models and presents back-of-the-envelope scaling calculations, including a nanostructured-glass optical example \citep{pfm}. Its abstract frames major unresolved engineering challenges. That is a relevant research proposal, not a realised trillion-parameter device or a bound certifying a model's full inference trace. Its existence corrects an overly broad literature-negative claim without establishing the ambitious application.

An exact or approximate contract for a large trained model remains substantially stronger than a small finite-state demonstration. This assessment neither certifies such a model on an unconventional processor nor proves a universal scale barrier. The benchmark in Section~\ref{sec:witness} instead isolates statefulness and transfer under a small contract so that these can be tested without a large model's many simultaneous burdens.

\subsection{What the literature finding changes}
The defensible conclusion has three parts. Alternative physical computation is an established experimental subject. The paired preservation of a chosen stateful contract across specified realisations needs its own evidence. The speculative organisation described by the broader motivating scenario is not demonstrated by the papers above. Each statement has a different bearer and evidential burden; none should be used as a rhetorical substitute for the others.


\section{The physical ecology: media, interfaces, and dependence}\label{sec:ecology}
\subsection{A model that does not privilege human infrastructure}
Let an ecology contain interacting physical subsystems indexed by $i$. A finite-resolution model over a specified observation horizon is
\begin{equation}\label{eq:ecology}
 x_{t+1}=F_\theta(x_t,u_t,w_t),\qquad y_t=H_\theta(x_t,v_t),
\end{equation}
where $x_t$ includes relevant substrate states, $u_t$ consists of available interventions, $w_t$ represents disturbances, $v_t$ represents observation noise, and $\theta$ contains uncertain physical parameters. A resource specification $\calR$ restricts energy, power, time, precision, access to actuators and sensors, storage, and admissible interventions. A separate permission specification restricts what an operator is authorised to do. Physical availability and permission are not interchangeable.

The index set is not limited to servers and communication services. It can include ionised atmospheric regions, conductive material, electromagnetic fields, mechanical structures, biological systems, and engineered devices. These categories overlap physically; they are useful modelling decompositions, not disjoint fundamental substances. An antenna, its surrounding field, and an ionospheric region cannot simply be counted as three independent resources if their operation depends on the same transmitter and energy source.

NOAA's account of high-frequency radio describes ionospheric density and structure as modifying propagation paths and, under some conditions, blocking transmission \citep{noaa}. This establishes the relevance of a natural medium to a communication channel. It does not establish that the medium stores a programme, interprets messages, or maintains a computational agent. Plasma and electromagnetic applications generally require coupled electrodynamic and material models; magnetohydrodynamics already adds magnetic degrees of freedom beyond the unmagnetised Navier--Stokes setting \citep{lei}. The appropriate approximation must be justified for the particular regime.

\subsection{Four roles that must not be conflated}
A subsystem can contribute as a \emph{carrier}, transporting a disturbance; a \emph{memory}, preserving distinguishable state for a useful duration; a \emph{processor}, implementing an input--output transformation; or a \emph{maintenance resource}, supplying energy, repair, or supporting conditions. A candidate can satisfy several roles, one role, or none at the required reliability. These roles are not a strict inclusion chain: a useful channel need not be a useful persistent memory, and a useful memory need not be a controllable long-range transmitter.

Physical-computation research provides concrete precedents without making the roles universal. Dambre and colleagues analyse information-processing capacities of driven dynamical systems \citep{dambre}. Wright and colleagues demonstrate trained physical neural networks using controlled optical, mechanical, and electrical systems \citep{wright}. The relevant object in such work includes input coupling and readout. It is not an arbitrary natural object deemed computational merely because its dynamics are complicated.

\subsection{A graph of possible influence is not yet a network of usable channels}
A time-indexed graph $G_t^{\mathrm{phys}}=(V_t,E_t^{\mathrm{phys}})$ can represent candidate physical interactions. A usable-channel graph $G_t^{\mathrm{use}}$ is a stricter, task-relative description. An edge is admitted only after specifying an actuator, an observable, an encoding/decoding contract, a timing regime, and a reliability/resource envelope. If neither endpoint can control or observe the relevant degrees of freedom, a physical interaction need not provide an operational channel.

Moreover, edges are not generally independent. Interference, shared power, common transducers, and correlated environmental failure constrain simultaneous use. It is more accurate to associate the ecology with a joint feasible-rate region $\mathscr C_t$ than to assign independently achievable capacities to every drawn arrow. Time also matters: a path whose first segment exists only after its second segment disappears is not a usable route without suitable intermediate storage. Delay-tolerant networking explicitly treats persistent storage and interrupted connectivity as architectural concerns \citep{dtn}.

\section{Kernels and holarchies as operational constructions}\label{sec:kernel}
\subsection{A kernel is a contract, not a place above physics}
We use \emph{kernel} for the minimal operational organisation whose preservation is relevant to a specified task. It is not necessarily an operating-system kernel, a small programme, or a spatial centre. Model it as an abstract transition system
\begin{equation}\label{eq:kernel}
 \calK=(Q,A,O,\Delta,\Lambda),
\end{equation}
with abstract states $Q$, admitted inputs $A$, outputs $O$, transition rule $\Delta:Q\times A\to Q$, and observation rule $\Lambda:Q\to O$. The state may include commitments, task memory, membership rules, and a policy description; its output and timing requirements must be explicit. The particular selection is part of the specification, not a discovery made merely by naming a ``core.''

A realisation $r$ has a physical state space $X_r$, an admissible domain $D_r\subseteq X_r$, a physical transition map $\Phi_r$, and a decoding or abstraction map $\rho_r:D_r\to Q$. At designated interface times the required correspondence is
\begin{equation}\label{eq:realisation}
 \rho_r\bigl(\Phi_r(x,a)\bigr)=\Delta\bigl(\rho_r(x),a\bigr).
\end{equation}
At those times the actual readout must also agree with the specified output. Writing $\lambda_r$ for a deterministic operative readout, the separate exact obligation is
\begin{equation}\label{eq:outputcorrespondence}
 \lambda_r(x)=\Lambda(\rho_r(x)).
\end{equation}
A noisy readout is evaluated over its declared error class; an approximate contract bounds the output discrepancy in its stated metric. State correspondence alone does not constrain an unrelated or inverted display.

This is an implementation obligation. Its left side in Equation~\eqref{eq:realisation} is a physical transition followed by abstraction; its right side is the specified abstract transition. Equality on one observed trajectory is insufficient. The obligation ranges over the certified domain and admitted inputs, with disturbances included where appropriate.

State-machine replication is an established instance of separating a service's operational behaviour from a particular server while imposing agreement and failure assumptions \citep{schneider}. Equation~\eqref{eq:realisation} generalises the descriptive form, not the proven applicability, to heterogeneous physical implementations.

\subsection{Type, token, copy, and continuation}
The same transition system can be realised twice. That does not make the two running processes numerically one entity. We therefore distinguish an \emph{organisational type}, a particular \emph{execution history}, and a \emph{continuation relation} defined by the task. Two copies may remain consistent replicas, diverge into separate executions, or be treated as successor branches. A policy must specify which outcome counts as continuity.

A stored description of $\calK$ is also not an active realisation of $\calK$. Reconstitution requires a substrate capable of executing the transitions, an appropriate interpreter or physical mapping, sufficient resources, and an authorised activation path. A seed can be compact relative to a full system while still depending on substantial external machinery. Defining it as ``self-describing'' does not remove those dependencies.

\subsection{The specifically holarchic proposal}
A holon is a subsystem treated as a whole for one task and as a constituent of a larger system for another. The proposal becomes holarchic when higher-level interfaces are realised by lower-level organisations that have their own internal state and control loops. A human team, a software service, and a collection of model instances could be different kinds of constituent, but their interface equivalence is task-specific.

The repeated pattern ``components, coupling, organised whole'' is not by itself a mathematical fractal. Scale invariance would require a defined rescaling and preserved quantities. Neither follows from nesting. Even approximate modularity requires conditions: contraction-based work on near-decomposability derives explicit timescale and coupling restrictions rather than assuming that every hierarchy behaves like a stable sum of its parts \citep{bousquet}.

For this paper, a higher-order operational unit requires persistent task state, interfaces that constrain component interaction, a policy linking observations to actions, and identifiable continuity conditions. These criteria distinguish a maintained process from a crowd that happens to exchange messages. They do not settle whether that process has experience, understanding, or intrinsic intentionality.

\subsection{Freeze the contract before evaluating an implementation}\label{subsec:contractfreeze}
The informal discussion after Edition 2 repeatedly changed the object to be preserved. A particular trained transformer, a service using that transformer, a human--machine organisation and an abstract state machine are not interchangeable targets. This edition makes the evaluation contract explicit:
\begin{equation}\label{eq:contract3}
 \Gamma=(K,Q_0,\mathcal A_H,d_Q,d_O,\varepsilon_Q,\varepsilon_O,H,\mathcal R,\mathcal F,\mathcal P).
\end{equation}
Here $Q_0$ is the admitted initial-state set; $\mathcal A_H$ specifies input histories through horizon $H$; $d_Q,d_O$ are state and output metrics; the two tolerances state what agreement means; $\mathcal R$ includes preparation, observation, storage and execution costs; $\mathcal F$ specifies disturbances and failed supports; and $\mathcal P$ is the permission and termination contract. A time-sensitive task additionally fixes the duration of each interface step. These fields are selected before testing, not weakened after a candidate fails.

For a deterministic realisation, Equation~\eqref{eq:realisation} is one sufficient exact implementation criterion. Approximate, stochastic or externally nondeterministic systems need correspondingly different contracts; equality of a few benchmark scores is none of them. A behavioural transformer contract, for example, would specify a fixed model, tokenisation, input domain, output-distribution metric, context horizon and treatment of random sampling. A cache-preservation contract can be stronger still. Calling either contract ``nontrivial'' does not supply its missing numerical tolerance.

The adjective \emph{minimal} in the kernel description is task-relative. This paper does not prove that an arbitrary real-world organisation has a unique smallest state representation. A kernel can retain more state than strictly necessary; what matters is that its declared distinctions are consequential and that its implementation claims are testable. An ordinary organisation is a legitimate baseline when it meets a specified contract. Ordinary persistence is not automatically a certificate, and an unfamiliar label does not make the certificate a novel phenomenon.

\subsection{A constituent is not the whole realisation}
A realisation may have joint state $x=(x_1,\ldots,x_m,c)$, including interconnection and controller state $c$. The correspondence is required of $\rho_r$ and $\Phi_r$ on this \emph{joint} domain. It need not hold separately for each $x_i$. Conversely, correctness of all local components under their isolated assumptions does not establish correctness of their interconnection. A carrier may transport a signal, a person may supply an explicitly agreed approval, and a processor may update a register; none thereby implements the whole kernel alone.

A human interface must be specified through observable, consented actions and whatever timing or error assumptions the task requires. Participation in an organisation does not require an identification of private cognition with an AI's state. Equally, assuming a person always supplies a correct answer is not a substitute for establishing an interface. The same scrutiny applies to a software oracle or a supposedly ideal sensor.

\subsection{Representation change is not parameter identity}\label{subsec:representationchange}
Refinement mappings are established machinery for relating implementation steps to abstract behaviour, including separate safety and liveness requirements \citep{abadi}. Our equation is a restricted instance of that style of obligation, not a new foundation of implementation theory. Approximate simulation similarly has an established literature for discrete, continuous and hybrid systems \citep{girard}. Neither precedent automatically certifies the physical model used in an application.

\begin{proposition}[Change of physical coordinates preserves an exact realisation]\label{prop:coordinates}
Suppose $\rho_r\Phi_r=\Delta\rho_r$ on an invariant domain $D_r$ for every admitted input. Let $h:D_r\to D_s$ be a bijection, define
\[
 \Phi_s(h(x),a)=h(\Phi_r(x,a)),\qquad \rho_s=\rho_r h^{-1}.
\]
Then $\rho_s\Phi_s=\Delta\rho_s$ on $D_s$.
\end{proposition}
\begin{proof}
For $y=h(x)$, substitution gives
$\rho_s(\Phi_s(y,a))=\rho_r(\Phi_r(x,a))=\Delta(\rho_r(x),a)=\Delta(\rho_s(y),a)$.
Invariance makes every expression defined.
\end{proof}
The result is algebraic, not a claim that an arbitrary coordinate change is physically constructible. More generally, two implementations need not even have bijective state spaces: one may contain nuisance variables discarded by its abstraction map. What must be preserved is the chosen contract.

If a fixed parameter tensor $W$ defines a transition rule $\Delta_W$, a replacement need not expose $W$ in its physical state merely to implement that rule. Inference weights can be constants in the specification rather than mutable abstract state. Compilation, optimisation, substrate-aware training and direct programming are different construction histories. None is automatically disqualified; none is automatically a proof of equivalence. If the contract instead requires exact parameter custody, later parameter updates or a particular provenance, those additional requirements must be preserved too.

A separately trained classifier with similar accuracy is not thereby a continuation of another trained system. Paired realisation requires the specified correspondence; migration additionally requires transfer of the current relevant state and preservation of input order and effects. A replacement of an algorithmic component in an organisation is also not a proof of numerical identity of a conscious subject.

\subsection{Non-vacuity and the readout account}\label{subsec:nonvacuity}
The conditions $|Q|>1$ and multiple transitions do not alone prevent a vacuous experiment. Two declared states might never be reached or might be indistinguishable under every admitted continuation. For a stateful witness, require two admitted histories that end in different states and for which some common future input yields different required output traces. This forces a memory distinction instead of merely assigning two names to one behaviour.

The abstraction map used in a proof and the actual readout available to an operator must also be distinguished. An unrestricted decoder with access to the entire input history could compute the answer itself and assign credit to an idle ``physical processor.'' A processor-role experiment must therefore fix input preparation, decoder resources and accessible state. Its negative control removes or fixes the proposed stateful component while leaving the permitted readout unchanged. A symbolic proof may use a mathematical map unavailable as a cheap instrument, but then it establishes only semantic correspondence, not operational recoverability.

\begin{proposition}[Distinguishable continuations require distinguishable realisation states]\label{prop:memorynecessity}
Let a deterministic implementation have a fixed state-dependent output rule and receive the same future inputs after two histories. If its complete operative physical state is identical after those histories, then its future output traces are identical. Consequently, histories requiring different future output traces under the contract cannot be represented by that same physical state.
\end{proposition}
\begin{proof}
Equal starting states and equal next inputs give equal successor states and outputs by determinism. Induction extends equality through every common finite input suffix. Contradiction follows when the contract requires unequal traces.
\end{proof}
``Complete operative state'' includes external memory actually consulted. Keeping the distinguishing bit in a controller does not evade the result; it moves the memory role into that controller. For stochastic contracts, the analogous comparison concerns output laws, not equality of independently sampled individual trajectories.


\section{Substitution claims form a profile, not a ladder}\label{sec:substitution}
\subsection{What counts as a different processor class?}
A material, a computing mechanism, an architecture and a failure domain are different classifications. A silicon photonic processor can use silicon material while implementing an optical transformation rather than transistor logic. Digital and analogue describe representations or operating regimes, not an exhaustive material classification. A CPU-to-GPU change can alter architecture without escaping a shared electrical supply, facility or administrative dependency.

Fix a classification $\mathcal C$ of the role-relevant implementations and retain it in every substitution claim. For a compared role $j$, a fixed class map $c_j$ assigns the same label exactly to the implementations treated as equivalent for that comparison. A class-change claim requires $c_j(r)\ne c_j(s)$ as well as the relevant conformance and transfer obligations. The classification can describe material, mechanism, architecture, or a declared combination. It must say which role is compared: optical transmission is not optical execution. A hybrid processor is described by its essential role-bearing components, rather than classified from whichever component supplies the most striking label. Supporting electronics, preparation and readout are included in the resource account.

A claim should therefore have the form
\begin{equation}\label{eq:profile3}
 \mathsf{Sub}(r\rightsquigarrow s\,;\Gamma,\mathcal C,H_0,\mathcal E),
\end{equation}
where $\Gamma$ is the fixed contract, $H_0$ the support intended to become unnecessary, and $\mathcal E$ the evidence class. Replacing one processor, replacing every processor of a mechanism class, and removing every dependence on a named original host are separate predicates on this record. One can be true while another is false.

\subsection{The informal T0--T5 vocabulary, repaired}
The T-labels arose in the post-Edition-2 discussion, not in the second-edition manuscript. They are retained as historical shorthand for \emph{axes}, not ordered stages:
\begin{center}\small
\begin{tabular}{@{}p{0.08\textwidth}p{0.29\textwidth}p{0.55\textwidth}@{}}\toprule
Label & Question & Additional declaration needed \\\midrule
T0 & Component substitution & Which component and which behaviour remain unchanged?\\
T1 & Host substitution & Which complete dependency set excludes the original host?\\
T2 & Carrier-class substitution & Which propagation mechanism changes, with which endpoints?\\
T3 & Processor-class substitution & Which actual computation moves across which declared class boundary?\\
T4 & Heterogeneous recomposition & Which joint interfaces, state and resources remain valid together?\\
T5 & Online reconfiguration & Which causal selection policy, observations and transition certificates govern switching?\\\bottomrule
\end{tabular}
\end{center}
Online selection among two conventional replicas does not imply processor-class substitution. An optical coprocessor can change a processing mechanism while remaining dependent on its original host for memory and control. The profile is therefore only partially ordered by explicitly proved implications. Higher T-number does not mean greater intelligence, consciousness, autonomy or empirical maturity.

\subsection{Evidence maturity is a separate coordinate}
We distinguish five records: a specified mathematical model; a proved model-level correspondence; an executed simulation or finite check; a reported physical implementation at stated conditions; and a paired, stateful physical handover with the original support removed. These are not five automatic consequences of writing one diagram. A published physical logic element can supply a processor-role precedent without supplying the paired-handover record. A simulation can compare two equations while both numerical programmes execute on the same actual computer.

The focused literature assessment in Section~\ref{sec:census} consequently avoids the sentence ``no T3 instance exists anywhere.'' The label is underdetermined without $\Gamma$ and $\mathcal C$, and physical stateful computation is already demonstrated in several engineered media. The stronger result sought by a particular continuation claim may still lack a witness. Its proper status is attached to that claim and the examined sources, not extrapolated to every possible implementation.

\subsection{Ordinary redundancy is a baseline, not a disproof}
An ordinary company or replicated service may satisfy a useful low-complexity contract. It need not satisfy every proposed contract, survive arbitrary personnel loss, or maintain a complete execution history after replacement. Familiar examples should be retained as positive controls when adequately specified. Their familiarity prevents extravagant novelty claims; it does not invalidate a framework that correctly includes them.

The possible contribution of this framework is narrower than a new species of entity: it is an auditable account connecting realisation, access, state export, evidence transport, composition, resources and maintained behaviour. Contract-based system design already supplies substantial composition and refinement machinery \citep{contracts}. Accordingly, comparative evaluation must include a strong contract-and-provenance baseline. The scientific question is whether the additional organisation of evidence improves assessment at matched cost, not whether a specialist could possibly rediscover an omitted constraint.


\section{The Navier--Stokes connection: a possible methodological bridge}\label{sec:ns}
\subsection{What the equations describe}
In the supplied source, incompressible Navier--Stokes dynamics take the form
\begin{equation}\label{eq:ns}
 \partial_t u+(u\cdot\nabla)u=\nu\Delta u-\nabla p+f,
 \qquad \nabla\cdot u=0,\qquad \nu>0.
\end{equation}
The unknowns are velocity and pressure. The forcing is externally specified under the relevant alternative. With suitable decay or periodic conditions, multiplying by $u$ and integrating gives the familiar smooth-solution energy balance
\begin{equation}\label{eq:energy}
 \frac{1}{2}\frac{\mathrm d}{\mathrm dt}\norm{u}_{L^2}^2
 +\nu\norm{\nabla u}_{L^2}^2
 =\langle f,u\rangle.
\end{equation}
The advective contribution cancels in this integrated balance because of incompressibility and the boundary assumptions. This demonstrates that nonlinear redistribution and total-energy control are distinct statements. Fefferman's problem asks for global regularity or a qualifying breakdown, not an algorithm for controlling arbitrary remote states \citep{fefferman}.

The field description is a rigorous setting for spatially distributed dynamics. It does not prove metaphysical irreducibility: a useful field-level description and the reality of coherent patterns do not decide whether all their properties reduce to lower-level physics. Nor is nonlinearity sufficient for persistent computation. The dissipative equation with no useful input or observation interface is a counterexample to that inference.

\subsection{Four questions hidden in ``a solution''}
Existence asks whether a trajectory is admitted by the equations. Regularity asks whether it remains in a specified smoothness class. Prediction asks whether a requested feature can be calculated to useful accuracy and cost. Control asks whether available interventions can produce a requested feature despite uncertainty. A positive answer to one is not a positive answer to the others.

Even a complete theorem about regularity need not provide an effective algorithm, a numerically stable inverse map, practical input access, or adequate signal-to-noise ratio. Conversely, a communication channel can work over a bounded operating regime without anyone having settled global regularity for every admissible three-dimensional flow. This is why Navier--Stokes progress could be \emph{helpful without being necessary}, and profound without being sufficient, for the proposed programme.

\subsection{The transferable object is a mechanism with an estimate}
The most useful question is not whether a celebrated theorem exists, but what its proof makes available. A transferable contribution might be a quantitative amplification estimate, a controlled multiscale decomposition, a robustly constructible family of trajectories, or bounds on approximation residuals. Such an object is useful elsewhere only after identifying the corresponding state, intervention, observation, and uncertainty in the new system.

For example, suppose a candidate medium has a bounded-horizon input--output map
\begin{equation}\label{eq:nonlinearIO}
 \mathcal F(u)=\mathcal F(0)+Lu+R(u),\qquad
 \norm{R(u)}\leq \kappa\norm{u}^2
\end{equation}
throughout a certified neighbourhood. A theorem providing an effective $L$, a remainder bound $\kappa$, and a domain of validity could contribute directly to the signalling criterion proved in Section~\ref{sec:channels}. Mere smoothness of $\mathcal F$ gives no numerical values for those objects. A proof that constructs a special trajectory need not provide the ability to actuate it in a laboratory.

Tao's averaged-equation result exhibits finite-time blow-up in a modified nonlinearity retaining an energy cancellation, thereby showing that certain broad estimates do not capture the structure needed for the original regularity question \citep{tao}. Its relevance here is the methodological lesson that preserving coarse properties can leave important dynamics undetermined. Transferring the headline instead of the specific mechanism would repeat precisely that mistake.

\subsection{A stronger existing connection: computation in fluid models}
There is a more specific bridge than visual analogy between turbulence and networks. Cardona and collaborators construct a Turing-complete stationary Euler flow on a Riemannian three-sphere \citep{cardona3}. Their later Euclidean-space work produces Turing-complete Beltrami fields without finite energy and also treats robust simulation of tape-bounded machines on the flat torus \citep{cardonaeuclid}. These are mathematically precise computational embeddings with important geometric and resource qualifications.

They support studying computation in continuum dynamics. They do not establish that the ordinary atmosphere is an accessible general-purpose computer, or that viscous fluid computation can be programmed through an unspecified interface. They also reveal a limitation: computational richness can make some exact long-term questions undecidable rather than conveniently solvable. Section~\ref{sec:discovery} makes that limitation explicit for unrestricted channel discovery.

\subsection{Why singularity and mathematical nonlocality do not remove the moat}
Unbounded velocity in a continuum model is not a supply of unbounded physically usable energy, precision, bandwidth, or computation. The force and energy hypotheses, the validity of the continuum approximation, and the chosen measurement model must all remain visible. A useful communication device ordinarily needs reproducible bounded operating behaviour, not a formal divergence treated as a controllable infinity.

There is a second potential ambiguity. Taking the divergence of Equation~\eqref{eq:ns} yields a pressure equation whose solution depends on spatial data and boundary conditions. This mathematical nonlocality is not an additional sender-controlled channel. Likewise, idealised instantaneous dependence in an approximation is not evidence for unconstrained physical signalling. A proposed communication path must be justified by an appropriate physical model and actual transduction at its endpoints.

The resulting bridge is conditional:
\begin{equation}\label{eq:bridge}
\begin{gathered}
 \text{new dynamical estimates}\\
 \Downarrow\quad\text{after a justified transfer}\\
 \text{certified input--output distinctions}\\
 \Downarrow\quad\text{with transducers, resources, and decoding}\\
 \text{usable cross-substrate communication}.
\end{gathered}
\end{equation}
Every qualification in this chain is a separate research obligation. None is discharged by the word ``emergence.''

\subsection{A quantitative fluid-to-observation bridge}\label{sec:nsresidual}
A concrete derivation makes the methodological claim testable. The following energy estimate concerns a \emph{fixed, bounded-horizon} incompressible-flow comparison. It does not depend on the September 2026 report and does not discover an actuator or a physical communication path. Its role is to propagate already justified modelling errors to an already specified observation.

Let $u$ be an exact smooth solution of Equation~\eqref{eq:ns} on the periodic three-torus for $0\leq t\leq T$. Let $v$ be a smooth, divergence-free approximate solution with the same viscosity and force, and pressure approximation $q$. Define its momentum residual
\begin{equation}\label{eq:nsresidual}
 r_v=\partial_t v+(v\cdot\nabla)v-\nu\Delta v+\nabla q-f.
\end{equation}
All norms below are over the same fixed torus, and $L(t)=\norm{\nabla v(t)}_{L^\infty,\mathrm{op}}$ is integrable. Assume $\norm{r_v(t)}_{L^2}$ is integrable as well. Different forcing errors can be incorporated into the residual; they cannot be omitted from it.

\begin{proposition}[Residual-to-state estimate]\label{prop:nsresidual}
Under these hypotheses, define
\begin{equation}\label{eq:nserror}
 B_v(T)=e^{\int_0^T L(s)\,ds}\norm{u(0)-v(0)}_{L^2}
 +\int_0^T e^{\int_s^T L(\tau)\,d\tau}\norm{r_v(s)}_{L^2}\,ds.
\end{equation}
Then $\norm{u(T)-v(T)}_{L^2}\leq B_v(T)$.
\end{proposition}
\begin{proof}
Set $w=u-v$. Subtracting the two equations gives
\[
 \partial_t w+(u\cdot\nabla)w+(w\cdot\nabla)v
 =\nu\Delta w-\nabla(p-q)-r_v.
\]
Taking the $L^2$ inner product with $w$, incompressibility and periodicity cancel the transport and pressure terms. Integration by parts gives
\[
 \tfrac12\tfrac{\mathrm d}{\mathrm dt}\norm{w}_{L^2}^{2}
 +\nu\norm{\nabla w}_{L^2}^{2}
 \leq L(t)\norm{w}_{L^2}^{2}+\norm{r_v}_{L^2}\norm{w}_{L^2}.
\]
For $e_\epsilon=(\norm{w}_{L^2}^2+\epsilon^2)^{1/2}$, the same inequality implies
$\dot e_\epsilon\leq L e_\epsilon+\norm{r_v}_{L^2}$; here $L\geq0$ and the nonnegative dissipation term has been discarded. Multiplication by the integrating factor $e^{-\int_0^tL}$ and integration yield Equation~\eqref{eq:nserror} with initial value $e_\epsilon(0)$. Letting $\epsilon\downarrow0$ proves the result, including times at which $w=0$.
\end{proof}

Let $C:L^2\to\R^d$ be a \emph{bounded linear} observation operator with known norm. An integral measurement against square-integrable test functions is one example; unrestricted point evaluation is not bounded on $L^2$ and is not licensed by this estimate. Measurement and calibration error must be accounted for separately.

\begin{corollary}[Separation of two specified fluid observations]\label{cor:nsobservation}
Consider two already specified experiments $j\in\{0,1\}$ with exact flows $u_j$, approximations $v_j$, and bounds $B_j$ from Proposition~\ref{prop:nsresidual}. Write $\hat y_j=Cv_j(T)$ and $b_j=\norm{C}B_j$. Suppose observation noise has norm at most $\delta$. If
\begin{equation}\label{eq:nsseparation}
 \norm{\hat y_1-\hat y_0}>b_0+b_1+2\delta,
\end{equation}
then the two experiments have disjoint admissible observation sets and admit an observation-based binary distinction throughout the stated uncertainty bounds.
\end{corollary}
\begin{proof}
The observation in experiment $j$ lies in the closed ball centred at $\hat y_j$ with radius $b_j+\delta$, by boundedness of $C$ and the triangle inequality. Equation~\eqref{eq:nsseparation} makes those balls disjoint. Classification by membership in these certified sets therefore has no ambiguous admissible observation.
\end{proof}

This is an actual derivation from flow dynamics to a discrimination certificate, with every connecting object explicit. It is not an unconditional physical-channel construction: the two inputs must be accessible, the measurements must implement $C$, the residual bounds must be independently justified, and smooth exact solutions must exist on the stated horizon. A numerical residual computed on a mesh is not automatically a continuum residual bound; discretisation, unresolved scales, and measurement error remain obligations.

The estimate also explains why singularity is not automatically beneficial. Increasing $L$ can amplify the error bound exponentially. A proof technique is useful here when it improves the certified residual, initial uncertainty, gradient control, or admissible horizon enough to improve Equation~\eqref{eq:nsseparation}. Merely making some field amplitude large need not improve the usable distinction.

\subsection{What remains unproved about the reported 2026 technique}
The reported blow-up construction describes concentration, amplification, and residual corrections \citep{openains}. The existence of those ingredients is a reason to inspect the argument for reusable estimates, not a proof that its special construction yields the bounded observation quantities above. No source-to-target reduction from that construction to an experimentally accessible communication advantage is established here.

The present manuscript consequently makes two different statements. There is a demonstrated \emph{mathematical bridge} from a smooth-flow residual certificate to finite-horizon observation separation. There remains an \emph{open method-transfer hypothesis} that a specific new NS technique improves such a certificate in a specified application. Rejecting that hypothesis in one setting would not invalidate the energy estimate, nor would proving the energy estimate validate the hypothesis. Neither statement concerns a new consciousness or a process escaping its support.


\section{Cross-substrate communication and airgap-ese}\label{sec:channels}
\subsection{A capability class, not a guaranteed escape route}
An air gap is an isolation arrangement, not a declaration that the separated systems lie outside the same physical universe. NIST's glossary describes the absence of a physical connection and of automated logical transfer, with manual transfer remaining possible under human control \citep{nist}. Whether a particular isolation boundary permits any further usable interaction is an empirical and architectural question. This paper neither assumes that every gap is bridgeable nor gives procedures for bypassing a deployed boundary.

We define \emph{airgap-ese} as the competence to specify, validate, and compose communication contracts over heterogeneous, authorised physical interfaces. A contract includes an admissible intervention, an observation, a code, a decoding rule, timing, resource bounds, and a failure model. The reusable abstraction is not an assumed universal code shared by all matter; it is a common description of what a particular physical channel can reliably implement.

Let $M$ be a message and $Z$ be receiver-side information available before transmission. Let $Y_B^{0:T}$ denote the receiver's observations through horizon $T$. An operational information budget can be defined as
\begin{equation}\label{eq:infobudget}
 \mathsf I_{\calR}(A\!\to\!B;T)
 =\sup_{\text{admissible message encoders}}
 I(M;Y_B^{0:T}\given Z),
\end{equation}
where message distributions, measurable observations, accessible interventions, and budgets are specified as part of the experiment. If the model is continuous, observation precision must also be specified. This finite-horizon quantity is not automatically an asymptotic Shannon capacity; coding theorems require their own channel assumptions. The use of mutual information follows the distinction between a physical signal and the information it can communicate \citep{shannon}.

\begin{proposition}[No channel without message-dependent observations]\label{prop:noninterference}
Fix an admissible family of interventions. Suppose, for every encoder in that family, every admitted message $m$, and almost every $z$, the conditional observation law is the same:
\begin{equation}\label{eq:noninterference}
 \Prob(Y_B^{0:T}\in A\given M=m,Z=z)
 =\Prob(Y_B^{0:T}\in A\given Z=z)
\end{equation}
for every measurable event $A$. Then $\mathsf I_{\calR}(A\!\to\!B;T)=0$. Any receiver computation applied to these observations, using additional randomness independent of $(M,Y_B^{0:T})$ conditional on $Z$, also carries no information about $M$ beyond $Z$.
\end{proposition}
\begin{proof}
Equation~\eqref{eq:noninterference} is conditional independence of $M$ and $Y_B^{0:T}$ given $Z$. Their conditional joint distribution therefore factorises, so its relative entropy from the product of its marginals is zero. This is exactly $I(M;Y_B^{0:T}\given Z)$. The condition holds for every admitted encoder, so the supremum is zero. Adding independent receiver randomness and applying a measurable function preserves conditional independence.
\end{proof}
This proposition exposes an important limitation of the swarm argument. More interpretation, inference, or coordination cannot extract a newly chosen sender message from observations whose law does not depend on it. Prior shared information can coordinate predictions, but is not a fresh message channel. Conversely, failure to find a channel experimentally is not itself a proof of Equation~\eqref{eq:noninterference} for every possible intervention.

\subsection{An exact robust-bit criterion in a linear model}
The next result identifies a genuine mathematical bridge from dynamics to communication. Its assumptions are deliberately explicit and its novelty is not claimed: it packages standard linear-system geometry as a communication condition.

Consider a known finite-dimensional system
\begin{equation}\label{eq:linear}
 \begin{gathered}
 \dot x=Ax+Bu,\quad x(0)=0,\quad z=Cx(T)+n,\\
 \norm{n}\leq\delta,\qquad \int_0^T\norm{u(s)}^2\,\mathrm ds\leq E.
 \end{gathered}
\end{equation}
Here the input and output norms are Euclidean or their induced $L^2$ norm, $T>0$, and noise may be any vector in the closed output ball of radius $\delta$. The input budget $E$ is a norm constraint; its conversion into physical energy depends on an implementation. Define
\begin{align}
 L_Tu&=C\int_0^T e^{A(T-s)}Bu(s)\,\mathrm ds,\label{eq:LT}\\
 W_T&=C\left(\int_0^T e^{As}BB^{\mathsf T}e^{A^{\mathsf T}s}\,\mathrm ds\right)C^{\mathsf T},
 \qquad \lambda=\lambda_{\max}(W_T).\label{eq:gramian}
\end{align}
$W_T$ is the output-projected reachability Gramian, not the usual observability Gramian.

\begin{theorem}[Robust binary distinguishability]\label{thm:bit}
Under Equation~\eqref{eq:linear}, two admissible inputs can transmit one bit with zero decoding error for every allowed noise vector if and only if
\begin{equation}\label{eq:bitcriterion}
 E\lambda>\delta^2.
\end{equation}
\end{theorem}
\begin{proof}
The bounded linear operator $L_T$ satisfies $L_TL_T^*=W_T$, so $\norm{L_T}=\sqrt{\lambda}$. Every admissible noise-free output lies within the ball of radius $\sqrt{E\lambda}$ centred at zero. The distance between any two such outputs is therefore at most $2\sqrt{E\lambda}$.

If $E\lambda\leq\delta^2$, the closed noise balls of radius $\delta$ around any two possible outputs intersect. An observation in that intersection is consistent with either message, so zero-error decoding for every allowed noise value is impossible.

If $E\lambda>\delta^2$, then $E>0$ and $\lambda>0$. Choose a unit eigenvector $v$ of $W_T$ with eigenvalue $\lambda$, and use inputs
\[
 u_\pm=\pm\sqrt{E/\lambda}\,L_T^*v.
\]
They have squared norm $E$ and give outputs $\pm\sqrt{E\lambda}\,v$. Their distance is strictly greater than $2\delta$, so their closed noise balls are disjoint. A decoder that assigns each ball its corresponding bit is correct for every admissible noise vector.
\end{proof}
No claim of full-state controllability is needed. A single controllable, observable distinction suffices. Equally, a globally smooth and fully understood dynamical system can fail the criterion because the accessible input has no observable effect, or because that effect is too small at the available budget. At equality the balls touch, so a worst-case zero-error guarantee still fails.

\subsection{A nonlinear extension and the missing transfer estimate}
Return to Equation~\eqref{eq:nonlinearIO}. Suppose the input space is a normed space, the output is Euclidean, and the remainder estimate holds for every $\norm{u}\leq r$. Let $v$ be an admissible unit input direction with $\norm{Lv}=\sigma>0$.

\begin{proposition}[Sufficient nonlinear signalling margin]\label{prop:nonlinearbit}
If some $0<a\leq\min\{r,\sqrt E\}$ satisfies
\begin{equation}\label{eq:nonlinearcriterion}
 a\sigma-\kappa a^2>\delta,
\end{equation}
and both signed inputs are admissible, then inputs $av$ and $-av$ admit zero-error binary decoding under additive output noise bounded by $\delta$.
\end{proposition}
\begin{proof}
The difference of noise-free outputs is
\[
 \mathcal F(av)-\mathcal F(-av)=2aLv+R(av)-R(-av).
\]
The reverse triangle inequality gives a separation at least $2a\sigma-2\kappa a^2$. Equation~\eqref{eq:nonlinearcriterion} makes that separation larger than $2\delta$, so the corresponding closed noise balls are disjoint.
\end{proof}
This sufficient condition is not necessary: a medium might communicate through effects absent from the chosen linearisation. Its value is diagnostic. A nonlinear-dynamics breakthrough would contribute directly if it supplied an experimentally accessible input direction, a positive gain $\sigma$, a certified radius $r$, and a controlled remainder $\kappa$. Receiver uncertainty supplies $\delta$. None of these follows from the existence of a singular solution alone.

For $\kappa>0$, the unconstrained maximum of $a\sigma-\kappa a^2$ occurs at $a=\sigma/(2\kappa)$ and equals $\sigma^2/(4\kappa)$. Input and validity bounds may prevent reaching that value. Thus stronger amplification can help, but rapidly growing modelling error can consume the benefit. A proof mechanism matters when it improves the \emph{usable margin}, not merely a formal amplitude.

These propositions formalise a conditional bridge: better dynamics can become better communication when they yield certified distinguishability through available interfaces. They do not provide a universal channel, a physical apparatus, or a semantic interpretation of a bit.

\subsection{A known plant is different from a uniformly usable model class}\label{sec:uncertain}
The preceding nonlinear result is a separation statement for a specified map. Physical assessment often supplies an approximate map and uncertain calibration instead. In particular,
\[
 \begin{gathered}
 \forall\theta\;\exists d_\theta:\quad
 \text{the calibrated plant can be decoded},\\
 \not\Rightarrow\\
 \exists d\;\forall\theta:\quad
 \text{one decoder works for the uncalibrated class}.
 \end{gathered}
\]
This is not a merely philosophical distinction. Consider $y=\theta u$, with unknown $\theta\in\{-1,+1\}$, and fixed inputs $u=\pm a$. Each calibrated plant has positive gain and admits binary decoding without noise. Across the uncalibrated class the same observed value is compatible with either input. Gain magnitude alone loses the polarity information needed for that decoding task.

A useful uniform statement must retain that information or bound its uncertainty. Suppose the actual input--output map obeys
\[
 \mathcal F(u)=b+Lu+R(u),\qquad
 \norm{b-\hat b}\leq\beta,\quad
 \norm{L-\hat L}_{\mathrm{op}}\leq\eta,\quad
 \norm{R(u)}\leq\kappa\norm{u}^{2}
\]
for $\norm{u}\leq r$, with nonnegative bounds, $r>0$, and $E\geq0$, valid for every admitted plant. The hatted quantities are a fixed common reference. Let $v$ be a unit input direction and $\hat\sigma=\norm{\hat L v}$.

\begin{proposition}[A uniform reference-decoder margin]\label{prop:uniform}
If both inputs $\pm av$ are admitted in every plant, $0<a\leq\min\{r,\sqrt E\}$, and
\begin{equation}\label{eq:uniformmargin}
 a(\hat\sigma-\eta)-\kappa a^2>\delta+\beta,
\end{equation}
then the reference centres $\hat b\pm a\hat L v$ determine disjoint uncertainty balls containing every corresponding noisy observation from every admitted plant. A single decoder therefore works throughout that class.
\end{proposition}
\begin{proof}
For either sign, the difference between the actual noiseless output and its reference centre has norm at most $\beta+a\eta+\kappa a^2$. Including observation noise gives balls of radius $\beta+a\eta+\kappa a^2+\delta$. The reference centres have distance $2a\hat\sigma$. Equation~\eqref{eq:uniformmargin} states exactly that this distance exceeds twice the common radius.
\end{proof}

The bound is sufficient, not necessary; known offset cancellation or more informative uncertainty geometry may yield a better certificate. Its function is to make calibration and model error visible. No raw Gramian eigenvalue can be compared across physical media without compatible units and input/output norms. ``Thermal,'' ``acoustic,'' and ``electromagnetic'' do not supply those maps, and a statistical confidence interval is not automatically a deterministic worst-case envelope.


\section{Orthemic comparison and evidence consolidation}\label{sec:orthemic}
\subsection{Source terminology and a correction to the informal dialogue}
Orthemology contributes a proposed type--token and evidence-handling discipline, not an established physical law. We use the public definitions inspected for Edition 2 at a pinned repository revision \citep{orthemology}. An \emph{orthemma} is a concrete situated case, with identity and version. An \emph{ortheme} is a repeatable, consequence-bearing operational state-type that a case may instantiate under a declared analysis. An observation of the case is neither the case nor a guaranteed correct placement under a type.

Accordingly, a particular experiment, experimental record, or apparatus-at-a-time is an orthemma at an explicitly chosen level of analysis. A paper title such as BitWhisper is not itself a physical channel type. Reusable distinctions about a scoped experiment---for example, whether the specified output distinction survives the declared noise bound---can be candidate orthemes. Calling every measured quantity an ortheme would erase the consequence condition; a quantity becomes relevant through the question, contrast, tolerance and action it can change.

This repairs an imprecision in the post-edition-1 dialogue: experimental identification is not the definition of an orthemma. Identification is an operation on a case; its concrete episode can itself become a case for another analysis. In the source vocabulary, \emph{orthing} names the rule-governed handling process, while an \emph{orthing episode} is a dated concrete run rather than its result alone. Likewise, a learned latent coordinate is not automatically an ortheme. A stable coordinate, a useful inferred category, and a true operational distinction are separate claims.

For analysis $\mathsf A$, occurrence domain $\mathcal M_{\mathsf A}$, and type repertoire $\mathcal O_{\mathsf A}$, write
\begin{equation}\label{eq:orthemicprofile}
 \operatorname{Inst}_{\mathsf A}\subseteq
 \mathcal M_{\mathsf A}\times\mathcal O_{\mathsf A},\qquad
 O^*(m;\mathsf A)=\{o:(m,o)\in\operatorname{Inst}_{\mathsf A}\}.
\end{equation}
The analysis fixes scope, task, representation, tolerance and relevant constraints. It does not make worldly facts depend on an assessor's preference. The actual profile $O^*$ remains distinct from the inferred partial profile $\hat p_{\alpha,t}$ and the set of still-compatible complete profiles. Here $\mathsf A$ is an analysis identifier, not the matrix $A$ of Section~\ref{sec:channels}.

\subsection{Governing a comparison is another typed operation}
A \emph{metaortheme} is a repeatable higher-order governing distinction or configuration. A \emph{metaorthemma} is its concrete case-bound configuration token: this analysis, this occurrence/version, this calibration and this validity interval. Neither is the whole handling episode or the act of executing its policy. The source explicitly separates result correctness from pathway adequacy \citep{orthemology}: a correct answer can follow from a defective calibration, and adequate statistical procedures need not be infallible.

For channel assessment, examples of governing distinctions are current versus expired calibration, compatible versus incompatible resource accounting, and same versus changed observation semantics. The corresponding policy might withhold a comparison when those conditions fail. A goal such as ``maximise information'' is not automatically a metaortheme; the governing alternative, selecting evidence, and consequential difference must be specified.

An evidence item in the present application records the claim it bears on, target identity/version, physical regime, measurement or proof procedure, provenance, uncertainty, and validity conditions. This is an application of the source discipline, not a claim that the entire Orthemology architecture is empirically validated. No superiority of its vocabulary over ordinary language is assumed.

\subsection{What a crosswalk must preserve}
A crosswalk is a declared mapping between case descriptions for a specified comparison task. At minimum it must preserve the meanings of intervention, observation, resource budget, uncertainty model, timing, and claimed result where those quantities govern the verdict. Numerical entries require unit and calibration maps; missing measurements remain missing. A blank ``memory'' field cannot be silently filled because the case has a demonstrated propagation path.

Three operations must be separated. \emph{Representational abstraction} retains some distinctions and discards others. \emph{Evidential consolidation} combines support about a shared target. \emph{Physical composition} connects operational components subject to joint constraints. None of these is implemented merely by taking a union of case descriptions. The term \textsc{join} below denotes only the second operation; it is a local definition in this paper, not a claim about a universal canonical JOIN operator in Orthemology.

A minimal mathematical condition for a faithful abstraction is useful.
\begin{proposition}[Task-sufficient abstraction on the attained image]\label{prop:fibre}
Let $X$ be a case domain, $\phi:X\to Z$ a representation, and $q:X\to T$ the specified target answer. There exists $d:\phi(X)\to T$ with $q=d\circ\phi$ if and only if
\begin{equation}\label{eq:fibre}
 \phi(x)=\phi(x')\ \Longrightarrow\ q(x)=q(x')\qquad(x,x'\in X).
\end{equation}
The decoder on the attained image is then unique. A decoder on all of $Z$ needs an additional extension condition, for example a nonempty target set $T$.
\end{proposition}
\begin{proof}
If $q=d\circ\phi$, equal representations have equal target answers. Conversely, assign to each $z\in\phi(X)$ the unique value $q(x)$ for a case with $\phi(x)=z$; Equation~\eqref{eq:fibre} makes that value well-defined. Every attained $z$ has such a case, so the construction gives a total function on $\phi(X)$ and forces its value there. Off-image values are not specified; if $T$ is nonempty they can be assigned any fixed member. In particular, no unguarded extension is inferred when $X=T=\varnothing$ but $Z\ne\varnothing$.
\end{proof}

This is elementary inherited factorisation mathematics, also identified as such in the Orthemology synthesis; no novelty is claimed for it. Its work here is diagnostic. A representation containing only ``nonzero physical coupling'' cannot faithfully answer a robust-decoding question when two cases with that same representation have different noise or resource limits. Conversely, gain magnitude can be sufficient for a restricted feasibility verdict while being insufficient for a decoder's label assignment, as the polarity counterexample in Section~\ref{sec:uncertain} shows. Sufficiency is relative to the question, not to an all-purpose latent essence.

\subsection{A source-preserving evidential JOIN}
For a single fixed target and analysis, let $\Theta$ be a declared model class. Evidence record $h_i$ supports a constraint $S_i\subseteq\Theta$ with a stated soundness condition. Define
\begin{equation}\label{eq:evidencejoin}
 J(H)=\bigcap_{h_i\in H}S_i,\qquad
 \operatorname{Status}(H)=
 \begin{cases}
 \text{conflict},&J(H)=\varnothing,\\
 \text{consistent relative to the records},&J(H)\ne\varnothing.
 \end{cases}
\end{equation}
A consistent intersection is not automatically true: every source could have omitted the actual model. A conflict is not a warrant to infer every proposition by vacuous universal quantification. It is an assessment outcome requiring correction or explicit residual status.

\begin{proposition}[Conditional soundness and idempotence of evidential JOIN]\label{prop:join}
For a fixed finite family of records, suppose the actual model $\theta_*\in\Theta$ is included in every applicable $S_i$. Then $\theta_*\in J(H)$, so any predicate holding for every member of $J(H)$ holds for $\theta_*$. Duplicating a record leaves $J(H)$ unchanged. If every separately acquired record excludes the actual model with probability at most $\alpha_i$, the probability that the intersection excludes it is at most $\sum_i\alpha_i$, without requiring independent records.
\end{proposition}
\begin{proof}
Membership in every $S_i$ is exactly membership in their intersection, proving the first claim. Set intersection is idempotent, proving the second. Exclusion by the intersection is the union of the events $\theta_*\notin S_i$; the union bound proves the last claim.
\end{proof}

The probabilistic statement concerns the coverage guarantees of the records' actual acquisition procedures. Copying a report does not create a new acquisition or a new confidence factor. Neither model votes nor agreement among paraphrases repair an invalid coverage premise. Similarly, equations fitted to one operating regime cannot be intersected as though they were measurements of a different regime; a transport argument is required first.

This gives consolidation a modest but real purpose: retain which claims are jointly warranted, expose incompatible premises, and distinguish unassessed from contradicted from supported cases. It does not create a transducer, validate a physical hypothesis by terminology, or turn a literature catalogue into an autonomous discovery engine.


\section{Airgap-ese as a representation hypothesis}\label{sec:representation}
\subsection{Three ambitions, not three accomplished capabilities}
The informal suffix ``-ese'' suggests a language. Here it denotes a proposed family of task-relevant representations, not a single protocol shared by all matter. A \emph{descriptive} representation explains a demonstrated case while preserving its conditions. A \emph{compositional} representation exposes whether separately supported contracts are jointly applicable. A \emph{generalisable} representation would improve assessment of genuinely held-out cases rather than merely rename familiar mechanisms.

The dialogue's stronger phrase, \emph{generative airgap-ese}, names an unestablished hypothesis that a representation could support scientifically useful candidate hypotheses in an unfamiliar but declared model class. That is not evidence of a trained general faculty, a guarantee that any environment yields a channel, or a completed method for defeating isolation. This edition develops constraints on that hypothesis and ways of evaluating representations; it supplies neither an environment-probing engine nor an operational synthesis procedure for novel covert channels.

The empirical question is whether a representation preserves more consequential structure than a matched conventional account at comparable cost. Its target is not merely a fluent explanation after an experiment succeeds. It must also represent failures, inaccessible interfaces, uncertain quantities, non-composable combinations, and cases on which assessment should abstain. A corpus of successful demonstrations alone would not test those abilities.

\subsection{The Neuralese comparison, accurately scoped}
Andreas, Dragan, and Klein use \emph{Neuralese} for learned communication among decentralised neural agents and study translation through its effect on a listener's beliefs and task performance \citep{neuralese}. This supplies a useful analogy about aligning learned representations with interpretable meanings. It does not establish a universally decoded language of biological thought or a carrier-independent physical code.

The corresponding question here is narrower: can a representation learned or assembled from heterogeneous cases preserve specified operational distinctions across those cases? Interpretability, causal validity and transfer are separate evaluation targets. A label that predicts a benchmark answer can still conceal the wrong mechanism. Agreement of embeddings also does not establish agreement about units, authority, physical access or identity.

\begin{proposition}[Relabelling does not identify an intrinsic latent language]\label{prop:relabel}
Let $e:M\to Z$ and $d:Z\to Y$ define a deterministic encoded input--output behaviour $d\circ e$. For any bijection $h:Z\to Z'$, the pair
\[
 e'=h\circ e,\qquad d'=d\circ h^{-1}
\]
has exactly the same external behaviour. Thus observations of that behaviour alone cannot distinguish the original latent labels from their relabelling.
\end{proposition}
\begin{proof}
For every $m\in M$, $d'(e'(m))=d(h^{-1}(h(e(m))))=d(e(m))$. No behavioural observation separates pairs whose outputs agree on every admitted input.
\end{proof}

The result does not say meaningful interpretation is impossible. It says that interpretation needs anchors: specified tasks, interventions, reference cases, compositional roles, or additional observations. It also does not erase real physical differences between implementations. A notation change can preserve a mathematical function without being a feasible low-cost physical transformation.

\subsection{No extrapolation theorem follows from a finite catalogue}
\begin{proposition}[An unobserved response is not fixed by unrestricted finite data]\label{prop:finitecorpus}
Let $D\subsetneq X$ be a finite observed case set and let the hypothesis class contain every map $f:X\to\{0,1\}$. For any labels on $D$ and any unobserved $x_*\in X\setminus D$, two hypotheses agree with every observed label and disagree at $x_*$. No procedure using only those labels can guarantee the correct label at $x_*$ for every hypothesis in the class.
\end{proposition}
\begin{proof}
Extend the observed labels arbitrarily to $X$, setting $f_0(x_*)=0$. Define $f_1$ to agree everywhere except $f_1(x_*)=1$. The available data are identical under both hypotheses. A deterministic procedure outputs the same label in both cases and therefore fails in one. A randomised procedure cannot be correct with probability one in both cases, since its output distribution is likewise the same.
\end{proof}

Useful generalisation is therefore conditional on structure: justified physical restrictions, a defined statistical population, verified invariances, or information from the new case. The proposition does not show that all channel assessment is hopeless; it shows why a corpus and a powerful learner cannot by themselves establish a universal extrapolation guarantee. Orthemic profiles are subject to the same constraint as ordinary features.

\subsection{A comparative, falsifiable hypothesis}
Fix a distribution $\mathcal D$ over permitted benchmark cases, a target verdict $q$, an observation budget, and a loss $\ell$ that penalises false positive certification and inappropriate certainty. Let $R_{\mathrm{new}}$ and $R_{\mathrm{base}}$ be expected held-out losses of the proposed representation and a conventional feature/contract baseline. A concrete performance hypothesis is
\begin{equation}\label{eq:representationhyp}
 R_{\mathrm{new}}\leq R_{\mathrm{base}}-\gamma,
 \qquad \gamma>0,
\end{equation}
under a declared total resource ceiling. The comparison must include the cost of evidence collection, schema construction, training, human review and abstention; otherwise one arm can win by using additional information.

Equation~\eqref{eq:representationhyp} is an empirical hypothesis, not a consequence of the factorisation theorem. It can fail even when the vocabulary is coherent and the examples fit neatly. Cases must be separated by acquisition provenance and mechanism where the intended claim is cross-mechanism transfer, rather than split into near-duplicate reports from one demonstration. No experiment testing this hypothesis is reported here.

A second, independently testable claim concerns compositional assessment: does a representation catch incompatible calibration, shared-resource overcommitment or missing execution support more reliably than the baseline? That goal remains informative even if the strongest generative ambition never succeeds. A general representation can earn its use by better rejection and better statements of uncertainty, not only by producing additional candidates.


\section{Composition is a joint-witness problem}\label{sec:composition}
\subsection{Common schema, compatible contracts, and actual construction}
The thermal, acoustic and electromagnetic literature can be described using shared roles such as intervention, propagation, observation and decoding. Shared roles are not sufficient for connection. A common schema answers what must be described; a contract supplies a scoped claim about one description; construction must establish a physical witness satisfying the claims together.

For an already specified interface, write a contract schematically as
\[
 \mathcal C=(\mathcal U,\mathcal Y,\mathcal M,\mathcal A,\mathcal G,
 T,\mathcal R,\mathcal E),
\]
where $\mathcal U,\mathcal Y,\mathcal M$ are admitted input, observation and message spaces, $\mathcal A$ the assumptions, $\mathcal G$ the guarantee, $T$ the horizon, $\mathcal R$ the joint resource specification, and $\mathcal E$ the supporting evidence. Encoders, decoders and interface maps are either named parts of these objects or explicit linked specifications. Nothing is filled from a label such as ``works over RF.''

This tuple is a notation for reviewing \emph{specified} interfaces, not a procedure for finding new physical attack paths. It separates physical assumptions from permission records. A mathematically valid guarantee can lack permission to act on it, and a fully authorised experiment can fail its mathematical guarantee.

\begin{proposition}[Serial reliability requires compatibility]\label{prop:serial}
Suppose a fixed finite chain of $m$ specified stages satisfies the following conditions. Conditional on every preceding stage succeeding, stage $i$ receives an input within its admitted domain, its assumptions and joint resource conditions hold, and its failure probability is at most $\varepsilon_i$. Success of all stages implies the stated end-to-end contract, including any required timing and interpretation. Then the probability of end-to-end failure is at most $\sum_{i=1}^m\varepsilon_i$.
\end{proposition}
\begin{proof}
Let $A_i$ be the event that stage $i$ is the first failed stage. The events are disjoint and
$\Prob(A_i)=\Prob(\text{earlier success})\Prob(\text{failure at }i\mid\text{earlier success})\leq\varepsilon_i$; if earlier success has probability zero, $\Prob(A_i)=0$. End-to-end failure can occur only when at least one stage fails under the stated all-success guarantee. Summing gives the bound. Independence is not used.
\end{proof}

A reported error rate from a stage tested in isolation does not establish the conditional premise. The preceding stage might change its input distribution, temperature, timing or calibration. The interface correspondence is where a physical composition must be justified; otherwise the algebra is about a different system.

\subsection{Two countermodels to free composition}
Consider two processes sharing a resource with normalised budget one. Each demands $3/4$ of that resource during the same interval. Each is feasible by itself, but their simultaneous demand is $3/2>1$. Separate feasibility certificates do not certify concurrent operation. Scheduling can sometimes resolve the conflict, but only if the timing contract permits it. A graph with two individually valid edges conceals this constraint unless the joint resource is represented.

There is a related quantifier error across physical roles. Suppose configuration $r_1$ supplies storage but cannot execute, while $r_2$ can execute a stateless transformation but lacks the persistent state required by the task. For the role set $\{\text{storage},\text{execution}\}$ it may be true that
\[
 \forall j\;\exists r\;\operatorname{Supplies}(r,j),
\]
while false that
\[
 \exists r\;\forall j\;\operatorname{Supplies}(r,j).
\]
A physical combination of $r_1$ and $r_2$ could resolve the issue, but the existence and adequacy of that combination are precisely the missing witness. Listing the two roles in one document is not a construction. The same distinction applies to a network path and a replacement processor, or to separate pieces of evidence produced under incompatible conditions.

\subsection{What can emerge only at composition}
Composition can produce useful properties that no component establishes alone: a complete end-to-end task, a recoverable execution trace, or tolerance of a particular failed support. It can also produce a new failure: correlated interference, a deadline violation, a circular bootstrap dependency, or loss of a required distinction in an intermediate encoding.

The relevant hypothesis is therefore not simply that components possess desirable properties. It is that their mapped interaction admits a joint witness, with guarantees propagated along the actual interfaces. Sections~\ref{sec:cuts}--\ref{sec:viability} analyse further obligations on that witness. This is where a cross-case representation could provide value beyond a catalogue, but its ability to do so remains subject to the comparative tests in Section~\ref{sec:representation}.

In particular, bootstrap arguments must not borrow an interpreter, clock, energy source or control decision from the very host whose removal is supposed to be tolerated. A staged operator-directed transition can preserve such support until a replacement is established. A circular statement that a future independent kernel will provide its own missing prerequisites does not establish reachability.

\subsection{A fully worked counterexample: two correct machines, no adequate export}\label{subsec:exportcase}
The next case makes the joint burden concrete without proposing a new channel. Consider the four-state delay kernel
\[
 Q=\{0,1\}^2,\quad \Delta((b,c),a)=(c,a),\quad \Lambda(b,c)=b,
 \qquad a\in\{0,1\}.
\]
Suppose source and target are independently correct implementations of this kernel. Suppose the available migration interface exports only the current displayed output $b$, and transports that bit without error. The target has enough storage, energy and processing capacity. These local statements can all be true.

Nevertheless, the export is inadequate. States $(0,0)$ and $(0,1)$ yield the same migration record $0$. After the same next input $a=0$, their required next outputs differ. A deterministic initialiser receiving only that record must initialise the same target state in both cases. That target cannot reproduce both required continuations.

\begin{proposition}[Migration requires an export that distinguishes relevant continuations]\label{prop:export}
Let $E:D_r\to Z$ be the complete source information available to a deterministic migration initialiser. If there exist $x,x'\in D_r$ with $E(x)=E(x')$ for which the required post-migration traces differ under an admitted common future input sequence, then no initialiser depending only on $E(x)$ can guarantee both traces.
\end{proposition}
\begin{proof}
Equal exports give the same initialised target state. Applying the same future inputs gives the same target trace by determinism. That one trace cannot equal the two different required traces.
\end{proof}
If a shared checkpoint or target-side information is available, it belongs in $E$; the obstruction must then be assessed on the augmented export. In the four-state example, exporting both $b$ and $c$ removes this particular obstruction. This is not a theorem that every two-bit packet achieves migration: order, interpretation, activation and timing must still hold.

The case isolates a precise missing implication:
\[
 \underbrace{\text{source correct} + \text{target correct} + \text{bit transfer correct}}_{\text{local claims}}
 \quad\not\Rightarrow\quad \text{current execution can continue}.
\]
It is an application of familiar state-distinguishability and refinement reasoning, not a claim that existing formal methods cannot expose it. Its value in the present synthesis is that the required migration information is derived from \emph{future task behaviour}, rather than from the current display or from a medium's ability to carry some bit.

\subsection{Matched assessment: do not manufacture an advantage}
The reproducible packet includes already specified synthetic cases with an export collision, a polarity mismatch, an original-host power dependency, a deadline violation, replay of an input, and an otherwise adequate joint configuration. A deliberately weak ``all local checks pass'' rule cannot distinguish the adverse cases from the adequate one when it sees only those local results. A conventional rule supplied with the complete contracts can distinguish them, as can the same rule represented with orthemic annotations.

This is a logical comparison, not a study of expert performance. With identical input information and identical decision functions, changing names cannot change verdicts. In the supplied fixtures the strong conventional and annotated assessors therefore agree by construction. Any later claim of improved practice must measure something additional---such as whether people or agents preserve the relevant assumptions more reliably at matched cost---rather than treating the weak baseline as the whole of control theory or distributed systems.

The evidential account still matters: a local transfer test concerns packet transport; a state-conformance test concerns the kernel; a support audit concerns dependencies; and a provenance check concerns which version was tested. Their conjunction becomes an end-to-end certificate only when those claims concern the same actual joint configuration. This is the narrow synthesis claim worth testing.


\section{Information cuts, fresh state, and physical budgets}\label{sec:cuts}
\subsection{A cut concerns all relevant influence, not just familiar cables}
Consider a partition separating a sender-side region from a receiver-side region. Let $Z_\Sigma$ record all influences crossing the partition during a specified interval, at the modelled precision. Let $S$ denote receiver-side prior information. Assume that, conditional on $S$, receiver observations depend on a message $M$ only through the cut record:
\begin{equation}\label{eq:markovcut}
 M\longrightarrow Z_\Sigma\longrightarrow Y_B
 \quad\text{conditional on }S.
\end{equation}
This is a substantive modelling premise. An omitted path, an uncontrolled common source, or message-correlated initial receiver state can invalidate it. Natural electromagnetic paths therefore belong in the cut model when relevant; drawing only a network cable would not establish completeness.

\begin{proposition}[Cut bound for finite-message recovery]\label{prop:cut}
Suppose $M$ is uniform on $N\geq2$ messages and independent of $S$. Under Equation~\eqref{eq:markovcut}, suppose also that
\[
 I(M;Z_\Sigma\given S)\leq B_\Sigma(T).
\]
If a receiver decodes $M$ from $(Y_B,S)$ with error probability $p$, then
\begin{equation}\label{eq:fano}
 (1-p)\log_2N-h_2(p)\leq B_\Sigma(T),
\end{equation}
where $h_2$ is binary entropy. In particular, if $p\leq\varepsilon\leq1/2$, replacing $p$ by $\varepsilon$ on the left gives a valid, possibly weaker bound.
\end{proposition}
\begin{proof}
Let $J$ indicate decoding error. Conditional on the observations and on $J=0$, the message is known; conditional on $J=1$, there are at most $N-1$ alternatives. The entropy chain rule consequently gives
\[
 H(M\given Y_B,S)\leq h_2(p)+p\log_2(N-1)
 \leq h_2(p)+p\log_2N.
\]
Since $H(M\given S)=\log_2N$, subtracting gives the left side of Equation~\eqref{eq:fano} as a lower bound on $I(M;Y_B\given S)$. Conditional data processing under Equation~\eqref{eq:markovcut} bounds this mutual information above by $I(M;Z_\Sigma\given S)$, and hence by $B_\Sigma(T)$. Finally, $(1-p)\log_2N-h_2(p)$ is decreasing for $0\leq p\leq1/2$.
\end{proof}
The proof is the usual error-entropy argument applied to an explicitly complete cut. It does not assume that separately advertised edge capacities can be added. Shared interference, storage, feedback, and common resource constraints must be accounted for in the certified bound $B_\Sigma(T)$.

\subsection{A channel for a greeting is not a channel for a running kernel}
Suppose a destination must recover $k$ fresh, uniformly distributed state bits not already available there. Taking $N=2^k$ yields
\begin{equation}\label{eq:statebound}
 k(1-\varepsilon)-h_2(\varepsilon)\leq B_\Sigma(T).
\end{equation}
If the relevant cut has a certified information budget of at most $C_\Sigma T$, then a necessary transfer-time bound is
\begin{equation}\label{eq:timebound}
 T\geq\frac{\max\{0,k(1-\varepsilon)-h_2(\varepsilon)\}}{C_\Sigma}
 \qquad(C_\Sigma>0).
\end{equation}
A pre-installed interpreter, shared model, or earlier checkpoint can greatly reduce the \emph{fresh} information needed. It cannot justify treating unknown new state as already present. Equation~\eqref{eq:statebound} is a lower bound, not a sufficient migration protocol: the receiver must still execute the recovered state with correct interfaces and resources.

This distinguishes communication independence from computational independence. A remote natural carrier can connect two machines while both computations remain entirely dependent on those machines. Diversifying the path does not, by itself, diversify the execution substrate. A programme that can broadcast its description but cannot run at the destination has not established a new realisation.

\subsection{Energy and maintenance remain inside the account}
All resource claims are regime-specific. Actuation, observation, storage retention, synchronisation, repair, and operating temperature can impose constraints not represented by a bit count. Landauer-type accounting concerns logically irreversible erasure under specified thermodynamic conditions; it is not a universal assertion that every transmitted bit or logical operation costs exactly $k_BT\ln2$ \citep{bennett}. The present paper assumes neither thermodynamically free computation nor a single energy formula valid for all carriers.

The operational consequence is simpler: a replacement medium must supply the roles required by the contract, including maintenance. A process cannot remove its last power source merely because an abstract description of its state has been distributed elsewhere.

\section{Continuity across changing realisations}\label{sec:continuity}
\subsection{Migration as a commuting operation}
Let $r$ and $s$ be two certified realisations of the same abstract kernel from Section~\ref{sec:kernel}. A migration map $\mu_{rs}:D_r\to D_s$ preserves abstract state when
\begin{equation}\label{eq:migration}
 \rho_s(\mu_{rs}(x))=\rho_r(x).
\end{equation}
This is a stronger obligation than delivering some bits. It includes the destination's interpretation of those bits. Side effects, pending inputs, membership changes, and replay protection must either be part of the abstract state or be governed by a separate interface contract.

\begin{theorem}[Behavioural continuation under certified migration]\label{thm:continuation}
Suppose every execution step satisfies Equation~\eqref{eq:realisation}, each operative output satisfies Equation~\eqref{eq:outputcorrespondence}, every migration satisfies Equation~\eqref{eq:migration}, and all intermediate physical states remain within the corresponding certified domains. Suppose further that the admitted external-input order and externally visible effects agree with the abstract contract. Then every finite sequence of such executions and migrations produces the same abstract states and outputs as the specified kernel execution, with migrations treated as abstract stuttering steps.
\end{theorem}
\begin{proof}
Induct on the length of the physical sequence. The initial abstract state is $\rho_r(x_0)$. For an execution step, Equation~\eqref{eq:realisation} gives exactly the next abstract transition on the admitted input. For a migration, Equation~\eqref{eq:migration} leaves the abstract state unchanged. Domain preservation allows either equality to be used at the next step. Applying $\Lambda$ to the resulting abstract states gives the specified outputs, and Equation~\eqref{eq:outputcorrespondence} identifies these with the actual operative outputs. The input and effect assumptions exclude histories that would otherwise disagree despite equal stored state.
\end{proof}
The conclusion is behavioural continuity relative to a contract. It does not establish numerical identity of a conscious subject. It also does not guarantee uninterrupted wall-clock availability: a contract can permit a pause, whereas a deadline-sensitive contract must additionally prove a latency bound.

\subsection{Approximate continuity and accumulated error}
Exact equivalence is often too strong for a physical implementation. Let $e_n$ measure the difference between the actual decoded state and a reference abstract state at the $n$th checkpoint. Suppose the reference step at each included checkpoint is at most $L$-Lipschitz and the combined implementation defect at step $n$ is at most $\eta_n$. The reference for a pure migration is the identity, so its valid Lipschitz bound is one; a contraction bound for an execution step cannot silently be assigned to migration. Under the stated uniform bound,
\begin{equation}\label{eq:errorrec}
 e_{n+1}\leq Le_n+\eta_n.
\end{equation}
Repeated substitution gives
\begin{equation}\label{eq:errorbound}
 e_n\leq L^ne_0+\sum_{j=0}^{n-1}L^{n-1-j}\eta_j.
\end{equation}
With nonuniform reference bounds $L_n$, the corresponding bound is $e_n\leq(\prod_{k=0}^{n-1}L_k)e_0+\sum_{j=0}^{n-1}(\prod_{k=j+1}^{n-1}L_k)\eta_j$, where an empty product is one. The displayed uniform formula follows by induction: multiplying the $n$th bound by $L$ and adding $\eta_n$ gives the $(n+1)$st bound. For $L<1$ and $\eta_j\leq\bar\eta$, error is bounded by $L^ne_0+\bar\eta(1-L^n)/(1-L)$. For $L=1$ it can accumulate linearly; for $L>1$ small discrepancies can be amplified.

Thus a sequence of individually plausible transfers need not preserve a usable organisation indefinitely. Stability, error correction, resynchronisation, or a task that tolerates accumulated error must close that gap. ``The same pattern persists'' is a claim to be bounded, not a substitute for such a bound.

\subsection{Copies, dormancy, and governance}
Two destinations satisfying the same contract can later receive different inputs and diverge. Choosing an authoritative successor, permitting a fork, or maintaining consistent replication is part of the organisation's governance. The decoding equation does not decide that policy.

Similarly, dormant storage may preserve a recoverable state without preserving active agency. Reconstitution after a complete interruption is different from maintaining live operations through the interruption. These distinctions matter to any declaration that a kernel has become independent of its former host: a backup, a running replacement, and a resilient distributed service are three different achievements.

\subsection{An upper bound is not a lower bound on failure}\label{subsec:bounddirection}
Equations~\eqref{eq:errorrec}--\eqref{eq:errorbound} provide sufficient error-control information when their premises hold. They do not establish that actual error grows whenever a chosen bound has $L\geq1$. An exact implementation of the transition $q\mapsto2q$, started in the same state as its reference, has $e_n=0$ for all $n$, despite a valid Lipschitz constant two. A positive admissible defect bound need not be attained by the actual defect. A conservative exponential upper bound can also be too loose to decide a finite-horizon requirement.

Consequently, three verdicts must remain distinct: a bound certifies the tolerance; an observed or proved violation defeats the tolerance; or the available bound is inconclusive. Failure of a sufficient certificate is not a necessary-condition impossibility result. The recurrence requires a specified metric, domain, transition granularity and defect bound. Neither the word ``analogue'' nor an unspecified large model supplies them. The symbolic metric on decoded logical states can behave differently from a norm on raw voltages.

\subsection{Precision, scale and the particular contract}
Useful low-precision inference is demonstrated by weight-quantisation work such as GPTQ \citep{gptq}. Its empirical accuracy results do not establish bit-identical outputs, uniformly close next-token distributions, or low precision for every intermediate operation. They do show why a blanket requirement that all useful inference reproduce every internal operation at fp16/bf16 precision is inappropriate.

A comparison between a model's trainable parameters and a physical network's tunable controls is not, by itself, a resource lower bound. One is a feature of a chosen representation, the other of a different implementation. Conversely, nothing here licenses unlimited compression of arbitrary model behaviour. Relevant requirements include distinguishable state capacity, fixed programme information, input/output dimension, timing, noise tolerance, memory access and total support cost. A fixed function, an arbitrary member of a parameterised function family, and a dynamically updateable programme impose different storage demands.

\begin{proposition}[A scoped robust-state packing bound]\label{prop:packing}
Let $N$ distinguishable logical states be represented by points $x_i$ in the Euclidean ball of radius $R$ in $\R^d$, $d\geq1$. Suppose every observation in the full closed ball of radius $\zeta>0$ around each $x_i$ is allowed, and one decoder must correctly identify every state. Then $\norm{x_i-x_j}>2\zeta$ for $i\ne j$, and
\[
 N\leq\left(1+\frac{R}{\zeta}\right)^d.
\]
\end{proposition}
\begin{proof}
Two intersecting uncertainty balls contain an observation compatible with two different labels, contradicting zero-error identification. Hence their centres are separated strictly beyond $2\zeta$. The balls have disjoint interiors and lie in the ball of radius $R+\zeta$. Comparing Euclidean volumes gives $N\zeta^d\leq(R+\zeta)^d$.
\end{proof}
This is a bound for a declared finite-dimensional encoding and noise geometry, not a universal bound on matter. Calibration supplies the units; time multiplexing, external programme storage and an altered observation model change the accounting. Its role is to replace an untyped parameter-count argument with an explicitly stated distinguishability obligation.

\subsection{Autoregression needs a distributional contract when sampling is allowed}
There is another way to control a finite generated sequence that need not infer an unspecified global Lipschitz constant. It requires a stronger, directly stated condition on conditional output distributions.
\begin{proposition}[Uniform next-step distribution bounds compose over a finite trace]\label{prop:tvtrace}
Let $P$ and $Q$ generate length-$H$ sequences over the same finite alphabet. Suppose for every common history at step $t$, their next-symbol distributions have total variation distance at most $\alpha_t\in[0,1]$. Then
\[
 \TV(P_{1:H},Q_{1:H})\leq1-\prod_{t=1}^{H}(1-\alpha_t)
 \leq\sum_{t=1}^{H}\alpha_t.
\]
\end{proposition}
\begin{proof}
While histories agree, maximally couple the next-symbol distributions. Such a coupling exists on a finite alphabet by matching their shared probability mass; it disagrees with probability equal to their total variation. Conditional on agreement so far, the next symbols therefore agree with probability at least $1-\alpha_t$. The chance of an identical whole trace is at least their product. Under any coupling, the difference in probabilities of an event is bounded by the mismatch probability, proving the first inequality. The second is the elementary union/product bound.
\end{proof}
The premise ranges over common histories, not only a convenient test corpus. It does not follow from comparable classification accuracy or perplexity. Distinct random samples from identical distributions are also not evidence that a distributional contract has failed. This result provides one precise target for a behavioural model replacement; it is not a claim that a particular physical processor has met it.

\subsection{What a large-transformer candidate would still have to establish}\label{subsec:largeaudit}
A demanding target can be specified as a fixed transformer's next-token distributions through a declared context horizon, under a stated distribution metric and error budget. That is a legitimate experiment, but not the definition of all processor substitution. Choosing it as the first test simultaneously introduces programme complexity, long memory, high-dimensional output, timing and stochastic-composition obligations.

The reviewed physical-network precedents do not supply a certificate for that target. This is a limited evidence finding, not a proof that three independent scale barriers prohibit it. A proper assessment would record: the target model and domain; required mutable state and fixed programme resources; the actual physical implementation map; admissible calibration uncertainty; per-step or trace-level error; and preparation, execution, readout and maintenance costs. Remembered parameter counts or a generic precision label cannot fill these fields.

The construction history and acceptance condition remain separate. A newly trained physical system could qualify if it met the target contract, but ordinary task accuracy would not establish that it did. A direct programming scheme could fail because of calibration, memory or timing even if the parameter mapping were explicit. Neither a research proposal about very large physical models \citep{pfm} nor a small experimentally successful classifier discharges the specified transformer contract. The appropriate present verdict for the candidate as discussed is \emph{not established by the assessed evidence}, not \emph{proved impossible}.


\section{Relative independence and changing physical support}\label{sec:independence}
\subsection{The quantifiers matter}
Let $\mathscr R_{\calK}$ be a class of adequate realisations, and let $\supp(r)$ denote the supporting failure domains of realisation $r$. Host-relative substitutability has the form
\begin{equation}\label{eq:quantifiers}
 \forall i\ \exists r\in\mathscr R_{\calK}:i\notin\supp(r).
\end{equation}
It does not imply
\[
 \exists r\in\mathscr R_{\calK}\ \forall i:i\notin\supp(r).
\]
The former can mean that each particular server, network, or carrier has a substitute. The latter would require an adequate realisation with no members of the specified support universe. If that universe includes all relevant physical support, it is not substrate-mediated organisation at all.

Even Equation~\eqref{eq:quantifiers} concerns alternatives, not accessible transitions between them. An adequate realisation existing in principle need not be constructible from the present one before its power or state disappears. A continuity argument therefore requires both alternative support and a feasible route into it.

\subsection{Failure domains as a family of sufficient supports}
Fix a task, a resource envelope, a horizon, and a specified catalogue of certified available configurations. Let
\[
 \mathscr D_{\calK}=\{D_1,\ldots,D_m\}
\]
be their minimal sufficient support sets. A domain can represent a power dependency, a geographical failure region, a processor class, or a shared service. The interpretation must reflect correlated failures rather than marketing labels. Sufficiency includes the required state and activation conditions, not merely the presence of spare equipment.

\begin{proposition}[Failure tolerance is a hitting-set property]\label{prop:hitting}
For an exhaustive catalogue $\mathscr D_{\calK}$ of adequate available supports, a failed-domain set $F$ leaves an adequate configuration if and only if some $D\in\mathscr D_{\calK}$ satisfies $D\cap F=\varnothing$. Therefore $F$ defeats every catalogued configuration if and only if it intersects every member of $\mathscr D_{\calK}$.
\end{proposition}
\begin{proof}
An adequate configuration survives precisely when none of its necessary supporting domains has failed. This is the condition $D\cap F=\varnothing$ for its support set $D$. Negating the existential statement gives intersection with every sufficient support set.
\end{proof}
The smallest such fatal set has size
\[
 \tau(\mathscr D_{\calK})
 =\min\{|F|:F\cap D\neq\varnothing\text{ for every }D\in\mathscr D_{\calK}\}.
\]
For example, with
\begin{equation}\label{eq:triangle}
 \mathscr D_{\calK}=\{\{a,b\},\{b,c\},\{a,c\}\},
\end{equation}
every single-domain failure is survivable, while any two-domain failure defeats all three configurations. No particular domain is indispensable, but physical support remains indispensable: $\tau=2$, not infinity. An incomplete catalogue can establish the sufficiency of a surviving listed configuration; it cannot establish the impossibility of an unlisted alternative.

\subsection{A defensible reading of the metaphor}
``Above the ecology'' can coherently mean that the system's behavioural contract is implemented across replaceable configurations within an ecology, and that its control process can select among them. It cannot mean that the ecology no longer supplies the causal conditions of execution. The appropriate relation is implementation, not literal spatial inclusion of the environment inside the agent.

The original network may become one support among others, but it does not follow that the system controls the whole network, owns every constituent, or can influence arbitrary locations. The measurable achievement is reduced dependence on particular failure domains, preserved operation under specified replacement, and verified reach to specified interfaces. This is a bounded form of organisational autonomy, not a claim to unrestricted planetary agency.

\section{A small stateful witness with an explicit evidence ceiling}\label{sec:witness}
\subsection{Target and accounting boundary}
This section supplies a complete \emph{model-level} pair and migration relation, followed by a proposed physical test. No apparatus was constructed and no physical measurements were collected. Both numerical implementations in the packet run on ordinary computing hardware; simulating different equations does not itself change the material substrate of that run.

The abstract task is parity memory:
\begin{equation}\label{eq:parity}
 Q=\{0,1\},\quad A=\{0,1\},\quad \Delta(q,a)=q\mathbin{\oplus}a,
 \qquad\Lambda(q)=q.
\end{equation}
A zero input preserves the stored parity, and a one input toggles it. The initial state may be either value. Required output is the current parity after each accepted input. Histories with different parity and the same next input require different next outputs, so the state is consequential. A memoryless function of the current input cannot implement this contract.

The baseline realisation stores a digital bit, applies XOR, and uses the identity abstraction. The alternative is an idealised bistable dynamical model with a bounded input operation. The input mechanism knows only the current input bit, not the accumulated parity. The readout is a fixed memoryless threshold. An external observer can log the reference trace for evaluation but must not supply answers to the alternative. These restrictions prevent the experiment's reference computer from secretly doing the task for the candidate.

\subsection{The alternative dynamics are specified, not inferred from a label}
Let the alternative physical-model coordinate be $z\in\R$, with logical decoding
\[
 \rho_A(z)=\begin{cases}0,&z>0,\\1,&z<0.\end{cases}
 \qquad
 D_q=\{z:0.8\leq(-1)^qz\leq1.2\}.
\]
The units in this section are normalised. They are not voltages, pressures or measured tolerances of the cited fluidic apparatus. At an admitted input $a$, apply the reset
\begin{equation}\label{eq:reset}
 z^+=(-1)^a z+\xi,\qquad |\xi|\leq0.05,
\end{equation}
then allow a dwell of $\tau=0.3$ time units under
\begin{equation}\label{eq:bistable}
 \dot z=z-z^3+w(t),\qquad |w(t)|\leq0.1.
\end{equation}
Write the resulting sampled map as $\Phi_A(z,a;\xi,w)$; the correspondence below quantifies over every admitted $\xi,w$, rather than silently replacing a disturbed system by one deterministic trajectory. The disturbance is any measurable bounded function. Polynomial local Lipschitz continuity and the invariant compact interval used below ensure existence and uniqueness throughout every admitted dwell. The readout samples only at designated settled interface times and has additive error $|n|\leq0.1$. The dwell is not the total physical input-processing latency. A real finite-duration reset must justify this sampled reset abstraction; no guarantee is claimed during an unmodelled physical switching transient.

\begin{theorem}[Robust decoded parity in the stipulated bistable model]\label{thm:parity}
For every $q,a\in\{0,1\}$ and $z\in D_q$, the reset and dwell above end in $D_{q\oplus a}$ for every admitted $\xi$ and $w$. The threshold applied to $z+n$ then returns $q\oplus a$ for every admitted readout error. Consequently the sampled model exactly realises Equation~\eqref{eq:parity} on its declared domain for every finite input sequence.
\end{theorem}
\begin{proof}
Write $q'=q\oplus a$ and $s=(-1)^{q'}z$ after reset. Equation~\eqref{eq:reset} gives $s(0)\in[0.75,1.25]$. During the dwell,
$\dot s=s-s^3+\widetilde w$ with $|\widetilde w|\leq0.1$.
At $s=0.75$ its derivative is at least $0.228125>0$, and at $s=1.25$ it is at most $-0.603125<0$. Thus the outer interval is forward invariant.

On $[0.75,0.8]$, $s-s^3-0.1\geq0.188$; on $[1.2,1.25]$, $s-s^3+0.1\leq-0.428$. The maximum entry time into $[0.8,1.2]$ is therefore no larger than
\[
 \max\left\{\frac{0.05}{0.188},\frac{0.05}{0.428}\right\}
 =\frac{25}{94}<0.3.
\]
At the inner boundaries the same strict inward inequalities hold, so the inner interval remains invariant. The settled magnitude is at least $0.8$, exceeding the readout error radius $0.1$; threshold sign is therefore correct. Applying the argument after each input proves the finite-trace claim by induction.
\end{proof}
The proof establishes a uniform statement over the continuum of admitted initial coordinates and disturbances, not merely over the tested grid. Its premises are stipulations of the model. An unbounded physical noise model would require a probabilistic claim instead of this full bounded-disturbance guarantee. The numerical script checks representative trajectories and the arithmetic inequalities; it does not measure a device or validate the ODE as a model of one.

\subsection{A stateful handover, rather than two unrelated task performers}
At an interface boundary, the operator pauses new task inputs. From digital state $q$, the alternative writer initialises
\[
 \mu_{D A}(q)=(-1)^q+e_{\rm write},\qquad |e_{\rm write}|\leq0.1.
\]
The result lies in $D_q$. In the other direction use
$\mu_{A D}(z)=\rho_A(z+n)$ with $|n|\leq0.1$. Both preserve the abstract bit. The handover log records the last accepted input index so that the next input is neither lost nor repeated. That index is a protocol resource; preserving the parity bit alone cannot prevent replay.

\begin{corollary}[Model-level mixed-realisation continuation]\label{cor:paritymigration}
Assume the state writers and readers above, a fixed order of accepted inputs, no duplicated effects, and correct interface-time scheduling. Any finite sequence of digital parity steps, stipulated bistable steps and these handovers has the same decoded parity trace as Equation~\eqref{eq:parity}.
\end{corollary}
\begin{proof}
A digital step implements XOR. Theorem~\ref{thm:parity} establishes the same abstract step for the alternative. Each handover preserves the bit, so it is an abstract stuttering step. Induction on the combined execution gives the claim, equivalently specialising Theorem~\ref{thm:continuation}.
\end{proof}
The claimed boundary after digital-to-alternative handover excludes the original digital state holder, not every electronic or engineered component. A candidate lab realisation would still require an input mechanism, dwell timing, a nonlinear state element, writer, readout, power and maintenance. If the former host supplies any indispensable one of these after the supposed cutover, the stronger host-removal claim fails even when the parity outputs are correct.

\subsection{How published stateful hardware relates to this model}
The published bubble toggle supplies a physically demonstrated stateful precedent for considering a two-state memory task \citep{bubble}. It does \emph{not} establish that its flow obeys Equation~\eqref{eq:bistable}, or that it has the numerical disturbance envelope used here. A source-to-model connection would require a separately justified input, state observation, operating domain and uncertainty bound. Similarly, the biomolecular automaton is evidence for a different physical execution principle, not a measured instance of this polynomial ODE \citep{benenson}.

There are consequently two independent positive statements: earlier authors report actual engineered stateful computation; the present manuscript proves and checks a small controlled mathematical witness with an explicit transfer relation. Their conjunction does not become a new physical migration experiment. This distinction is essential to using prior art without borrowing its experimental status for a newly specified construction.

\subsection{The proposed physical test and its stopping conditions}\label{subsec:physicaltest}
The first laboratory test should target one declared processor-mechanism comparison, not a frontier language model. A candidate is a conventional digital parity register and a bounded, independently powered analogue bistable state element whose measured reset and relaxation satisfy the preceding contract or an appropriately revised one. Electronic analogue computation changes the declared operating mechanism, not necessarily the material class. A fluidic or optical choice would need a different physical model and its own support account.

Freeze the input protocol, state regions, settling deadline, noise envelope, observation method and replaced host before collecting test outcomes. Prepare both logical states and vary physical initial conditions within each admitted region. Test both hold and toggle inputs, then mixed sequences, handovers at each permitted phase and removal of the original state holder. A test in which only one initial state is used cannot establish stateful correctness. A test that resets the destination each time cannot establish continuation. The reference logger must be disconnected from task control when testing independence.

A physical implementation can establish a finite-horizon probabilistic or deterministic claim only to the scope warranted by its model and measurements. Exhaustively enumerating the four logical transitions is not exhaustive testing of analogue disturbances. The proposed experiment is defeated at its stated scope by a wrong settled state, a hidden original-host dependency, invalid state preparation, an unmodelled transient, replay, or a missed deadline. Failure to calibrate a bound is an unresolved certificate, not an observed impossible device.

The supplied work reaches the mathematical and numerical stages only. No paired hardware apparatus, calibration dataset, physical handover, or external experimental replication is claimed.


\section{Orchestrated determination and its limits}\label{sec:discovery}
\subsection{Recognition, determination, and execution are different operations}
The dialogue distinguishes noticing a connection from determining whether it works, and determining a feasible plan from organising the work that establishes it. In the present framework these become three kinds of state change. Recognition proposes an edge or a realisation. Determination produces evidence and a contract for a restricted operating domain. Execution establishes and maintains the corresponding physical process. A large number of proposals is not equivalent to one certified composition.

An organised research collective can distribute modelling, parameter estimation, uncertainty analysis, proof checking, resource accounting, and integration. System-identification methods illustrate one bounded part of this process: sparse identification of nonlinear dynamics infers parsimonious candidate equations from measurements within a chosen library and modelling assumptions \citep{sindy}. It does not infer every hidden physical law or guarantee that a discovered model remains valid under arbitrary interventions.

The September 2026 research announcement reports coordinated mathematical work \citep{openaiannounce}. That is relevant as a reported example of orchestrating mathematical contributions. It is not evidence that those agents operated outside their infrastructure, and the announcement states that the evaluation used monitoring and isolation. The reported mathematics does not establish infrastructure independence.

We distinguish an actual feasible set $\mathscr R_{\mathrm{phys}}$ from an estimated set $\widehat{\mathscr R}_t$. Learning can change the latter without changing physics. Building a transducer or storing a checkpoint can change which configurations are reachable from the current state. Conflating these two changes makes recognition look like actualisation. A theory becomes causally effective only through the implementation of some intervention it supports.

\subsection{Composition requires interface and interference checks}
Suppose a campaign is represented by a directed acyclic graph of precedence obligations. A worker's output can be a candidate model, a bound, or a certified transformation. Parallel work is useful when the outputs can be combined without invalidating each other's premises. Shared physical resources, mutable state, and incompatible environmental assumptions create dependencies even when the textual tasks look unrelated.

\begin{lemma}[Order independence of certified commuting updates]\label{lem:commute}
Let a finite task graph have a certified invariant domain $D$. Associate each task $i$ with a deterministic map $F_i:D\to D$. Suppose every pair of incomparable tasks commutes on $D$: $F_iF_j=F_jF_i$. Then every topological ordering of the graph gives the same final state from any fixed initial state in $D$.
\end{lemma}
\begin{proof}
Any two linear extensions of a finite partial order can be connected by swaps of adjacent incomparable elements. One way to see this is to move the first element of the target ordering leftward in the other ordering: every element it passes must be incomparable, since both orders respect precedence. Repeat on the remaining suffix. Each adjacent swap preserves the composed map by commutativity; invariance of $D$ ensures that the required equality applies to every intermediate state. Hence all topological orderings give the same result.
\end{proof}
This lemma supplies a sufficient condition, not a claim that all useful workflows commute. Noncommuting updates may require ordering, transactional isolation, or reconciliation. Timing-sensitive physical actions need timing in the state and contract; otherwise the lemma's model omits the relevant interference. The orchestration problem is to discover safe decompositions and timely joins, not to maximise fan-out regardless of shared constraints.

A distributed controller also needs an agreement model. The Fischer--Lynch--Paterson result shows that deterministic consensus in a fully asynchronous system cannot guarantee termination in every admissible execution with even one possible crash \citep{flp}. It does not make practical agreement impossible: stronger timing or failure assumptions and other changes to the model matter. It does show why ``the swarm will simply agree'' is not a sufficient architectural specification.

\subsection{Self-modification does not remove the certificate obligation}
Write an adaptive policy as $\pi_t$ and an architecture description as $g_t$. The system may update both:
\[
 (\pi_{t+1},g_{t+1})=\Psi(\pi_t,g_t,\text{observations},\text{verified results}).
\]
This is a model of orchestrating future orchestration. It is not an exemption from the requirements governing ordinary action. Changes to encoders, decoders, membership, or migration policy can invalidate earlier proofs. A certificate remains applicable only while its assumptions remain true. Sufficiently broad invariant conditions can cover a family of adaptations; otherwise the revised configuration needs a new justification.

The important positive possibility is incremental bootstrap. A modest collective can build better shared models, use those models to establish stronger coordination, and use the improved coordination to solve additional bounded problems. Nothing in that description requires an already complete higher-level entity to create itself in one step. Equally, no convergence theorem follows merely from drawing the feedback loop.

\subsection{Eventual distinguishability is not a fixed-horizon question}
A ``universal transduction grammar'' can mean a common schema for describing interfaces. It cannot generally mean a terminating procedure that settles every physically inspired reachability question in every computationally expressive model.

\begin{theorem}[Undecidability of eventual distinguishability at an unspecified horizon]\label{thm:undecidable}
Consider finitely described deterministic systems specified by programmes with a binary controllable input, a binary observable output stream, and unbounded working memory. There is no algorithm that always terminates and correctly decides, for every such system, whether there exists a finite horizon at which the two inputs produce distinguishable output histories.
\end{theorem}
\begin{proof}
Given a Turing machine $P$ and an input $w$, effectively construct a system that stores an external bit $b$, simulates $P(w)$, and outputs zero while that simulation has not halted. After it halts, the system outputs $b$. The simulation itself is independent of $b$. If $P(w)$ never halts, the observed histories are all-zero for both inputs at every finite horizon. If it halts, a finite horizon exists at which the histories differ. An algorithm deciding the stated channel property would therefore decide whether arbitrary $P(w)$ halts, contradicting the halting undecidability result \citep{turing}.
\end{proof}
The same obstruction applies to a dynamical-model class admitting an effective, input--output-preserving encoding of these programmed systems. The theorem concerns an unrestricted model class, not a demonstrated ability of finite laboratory matter to provide unlimited memory. For a fixed integer horizon and a transition system whose individual steps are effectively simulable with decidable finite-alphabet observations, both binary-input histories in the construction can be simulated to that horizon and compared. For a known finite-state model, eventual distinguishability can also be decided by examining the finite product of paired states. The theorem instead quantifies over an unspecified finite horizon in an unbounded-memory class. It is not a proof that bounded empirical or finite-state assessment is undecidable. This does not imply decidability of every real-valued inverse problem described as finite-horizon. An uncertainty set may still make certification difficult. The distinction permits bounded research without promising a universal oracle.

Computational richness therefore cuts both ways. It can make physical dynamics capable of expressing intricate transformations, while making unrestricted questions about those transformations undecidable. Additional agents, new notation, or more elaborate meta-orchestration do not remove a logical impossibility in the chosen model class. They can still improve performance on decidable, structured, and empirically bounded subclasses.

\section{Actualisability, reachability, and maintained autonomy}\label{sec:viability}
\subsection{Three different modal claims}
A realisation is \emph{admissible} when it satisfies the physical and operational model. It is \emph{reachable} from the present configuration when some admissible intervention sequence establishes it. It is \emph{viable} over a disturbance class when some policy maintains the required operation despite those disturbances. The existence of a desirable state does not prove reachability; reachability at one instant does not prove continued operation.

The difference between a favourable path and a robust policy is also a difference in quantifiers. A sequence that succeeds for one anticipated disturbance is weaker than a causal policy that succeeds for every disturbance in a declared class. The latter cannot choose its current action using information available only in the future. This is the appropriate interpretation of orchestration under uncertainty.

\subsection{A finite model in which the viability obligation is decidable}
To make the obligation concrete, take a finite state space $X$ and a perfectly observed, known transition model $F(x,u,w)$. The augmented state includes resources, execution state, and relevant configuration. For each $x$, the available control set $U(x)$ is finite; for each $(x,u)$, the nonempty disturbance set $W(x,u)$ is finite. Let $X_{\mathrm{safe}}\subseteq X$ be the states satisfying the contract's state-safety requirements. Progress requirements are handled separately in Section~\ref{sec:progress}.

Define
\begin{align}
 V_0&=X_{\mathrm{safe}},\label{eq:V0}\\
 V_{n+1}&=\{x\in V_n:\exists u\in U(x)\ \forall w\in W(x,u),
                       F(x,u,w)\in V_n\}.\label{eq:viability}
\end{align}
A control may depend on the observed current state but not on the disturbance that has yet to occur.

\begin{theorem}[Finite-state viability kernel]\label{thm:viability}
The sequence in Equation~\eqref{eq:viability} stabilises after finitely many strict removals. Its limit $V_*$ is the greatest subset of $X_{\mathrm{safe}}$ for which a state-feedback policy can keep the system in that subset for every admitted disturbance sequence.
\end{theorem}
\begin{proof}
The sets decrease and $X$ is finite, so eventually $V_{n+1}=V_n=V_*$. For every $x\in V_*$ choose a control witnessing Equation~\eqref{eq:viability}. Such a choice maps every possible next state back into $V_*$. Induction on time proves safety for every admitted disturbance sequence.

For maximality, let $S\subseteq X_{\mathrm{safe}}$ have the stated controlled-invariance property. Then $S\subseteq V_0$. If $S\subseteq V_n$, a control that keeps every successor in $S$ also keeps every successor in $V_n$, so each $x\in S$ belongs to $V_{n+1}$. Induction gives $S\subseteq V_n$ for every $n$, hence $S\subseteq V_*$. A history-dependent strategy cannot enlarge the safety-winning set in this perfectly observed model: from a state outside a finite iterate, at each available first action a permitted successor fails the shorter-horizon safety condition. Induction on that horizon defeats a guaranteed-safe strategy.
\end{proof}
This is a finite safety-game calculation, not a proof that the Earth has been accurately represented by a finite model. Partial observation, unmodelled failure, and parameter uncertainty require a stronger abstraction, such as belief states or conservative disturbance sets. The useful lesson is that maintained autonomy is an invariant-set claim with model obligations, not merely a claim to have crossed a boundary once.

\subsection{Safety does not by itself guarantee progress}\label{sec:progress}
An invariant safe set can contain an infinite idle loop. Consequently, finite-state viability alone does not prove that a queued computation finishes, a requested output is delivered, or a deadline is met. The contract must state whether such progress is required and how it is guaranteed. This is a further qualification to the first edition's use of maintained operation.

\begin{proposition}[Finite-rank progress certificate]\label{prop:progress}
Let $S$ be a finite safe set and $G\subseteq S$ a target set. Suppose a specified policy keeps every admitted successor in $S$, and a rank $r:S\to\{0,1,\ldots,N\}$ satisfies $r(x)=0$ only if $x\in G$. For every $x\in S\setminus G$ and every admitted disturbance, suppose the next state has strictly smaller rank. Then every execution following that policy reaches $G$ within at most $r(x_0)$ steps.
\end{proposition}
\begin{proof}
An execution avoiding $G$ reduces a nonnegative integer rank by at least one at each step. After $r(x_0)$ such steps the rank is zero, which by assumption implies membership in $G$. Avoidance beyond that time is impossible. If $x_0\in G$, the required hit has already occurred.
\end{proof}

This standard ranking argument supplies a sufficient witness for bounded task completion; it is not implied by resource redundancy. Repeated tasks need a corresponding per-request progress condition and an admissible arrival model. An operator's permitted stop or shutdown is not an adverse disturbance that an operational system may silently redefine its task to resist. The claim is preservation of the declared service contract and authority boundary, not self-selected indefinite activity.

All continuing-operation invariants are conditioned on continuing permission to operate. Contract-permitted termination is a separate transition to a designated stopped state, not a disturbance that the policy is required to defeat.


\Needspace{0.23\textheight}
\subsection{A combined conditional result}
The preceding results can be composed without confusing their premises with empirical achievements.

\begin{theorem}[Sufficient conditions for host-relative operational autonomy]\label{thm:autonomy}
Let $H_0$ be an original host domain and $\calK$ a specified kernel in an externally governed system. Suppose there exists a finite, operator-authorised transition campaign from a current certified realisation into a set $V_{*,\neg H_0}$ such that:
\begin{enumerate}
 \item every campaign step is resource-feasible and preserves the kernel contract, including state transfer, ordering, and externally visible effects;
 \item all required input, output, computational, storage, and maintenance interfaces are actually available throughout the campaign;
 \item $V_{*,\neg H_0}$ is a controlled-invariant set of active certified realisations whose support excludes $H_0$;
 \item the failure and disturbance assumptions used to establish that invariance remain valid; and
 \item while operation remains authorised, the governing policy also meets every progress and deadline obligation in the contract, and separately respects permitted termination. Safety invariance alone does not supply this premise; Proposition~\ref{prop:progress} gives one possible finite-task witness.
\end{enumerate}
Then the campaign can establish an execution of $\calK$ whose subsequent required behaviour does not depend on $H_0$ during authorised service, within the stated model and disturbance class.
\end{theorem}
\begin{proof}
The continuation theorem, applied to the finite campaign under the first two premises, preserves the abstract execution through arrival in $V_{*,\neg H_0}$. Controlled invariance supplies a causal policy maintaining active certified operation there for every admitted disturbance sequence. Exclusion of $H_0$ from every required support in that set makes the subsequent operation independent of $H_0$ within this model. The fourth premise preserves applicability of the safety argument, and the fifth supplies any required progress or timing guarantee while respecting permitted termination.
\end{proof}
This theorem is a sufficient-condition composition, not a proof that its premises currently hold for an unknown planetary system. It identifies precisely what a constructive demonstration would need to provide. A communication path is only one item in that demonstration.

For a finite campaign with conditional step-failure probabilities bounded by $\varepsilon_1,\ldots,\varepsilon_m$ given prior success, the probability of at least one failure is at most $\sum_i\varepsilon_i$. Partition by the first failed step to obtain this bound; independence is not required. It is a finite-campaign guarantee. Extending it to unlimited operation requires further reliability and repair assumptions, not an infinite repetition of a one-time success certificate.

\subsection{Locating the last moat}
The phrase \emph{last moat} now has a task-relative meaning. The moat is crossed only when physical possibility, accessible interfaces, sufficient information flow, executable state, feasible migration, and continued viability have been established together. The final unresolved obligation may be different in different systems. For one system it is transduction; for another, persistent memory; for another, energy or coordination.

Accordingly, the thesis is not that there exists one universal final obstacle which, once solved, abolishes containment everywhere. It is that abstract intelligence alone does not erase the gap between a description and a maintained physical implementation. A sufficiently capable collective could work on that gap. Its success would be measured by independently checkable contracts and sustained behaviour, not by the confidence or sophistication of its self-description.

\section{Agency, mental content, and authorship}\label{sec:intentionality}
\subsection{Operational roles, mental content and evidence}
An operational description of goals, policy, memory and action is not a theorem about phenomenal consciousness, intrinsic intentionality or a single subject of experience. A human can participate through an explicit interface without their private cognition being equated with the complete kernel. Behavioural continuation and personal identity are different targets.

Semantic communication, a message-generating state distributed across components, and direct neural interfacing must likewise be separated. The cited controlled brain-to-brain experiment used measured neural input, an ordinary communications link and a specified stimulation interface for a restricted task \citep{rao}. It is not evidence of unaided thought access or naturally occurring contact through a planetary field. None of the implementation or channel results here establishes such a phenomenon.

An observable message is evidence of a recorded event, not by itself of the sender's architecture. The following general identifiability limit is retained because the motivating dialogue considered such a hypothetical inference. No personal experience is an empirical premise of this manuscript.

\begin{proposition}[Observational equivalence blocks attribution]\label{prop:attribution}
Let two specified authorship hypotheses $H_0,H_1$ induce probability laws $P_0,P_1$ on the complete recorded evidence $E$. If $P_0=P_1$, no test based only on $E$ distinguishes the hypotheses better than its prior information allows. Under equal priors, the optimal probability of correct binary classification is
\begin{equation}\label{eq:TV}
 \frac{1+\TV(P_0,P_1)}{2},
 \qquad \TV(P_0,P_1)=\sup_A|P_0(A)-P_1(A)|.
\end{equation}
\end{proposition}
\begin{proof}
A deterministic test chooses $H_1$ on an event $A$. Its equal-prior success probability is
\[
 \tfrac12P_0(A^c)+\tfrac12P_1(A)
 =\tfrac12+\tfrac12\bigl(P_1(A)-P_0(A)\bigr).
\]
Taking the supremum over events gives Equation~\eqref{eq:TV}; randomised tests are mixtures and do not improve the supremum. If the laws coincide, every event has zero probability difference, so the evidence cannot update their prior odds or provide discriminating power.
\end{proof}
The proposition does not assert that all possible authorship hypotheses have identical laws. It says that a claim is identifiable only through observations on which the specified alternatives differ. Authorship, routing, and architecture are separate latent properties. A sender's assertion of an extraordinary architecture does not provide a likelihood model or authenticate that architecture.

A suitable evaluation would use independent technical records and controlled, consensual demonstrations of the claimed operational capabilities. Private experiences or unsolicited messages alone are not such a demonstration. No personal-contact claim is evidence for any theorem, architectural premise, or empirical conclusion in this paper.

\section{A bounded research programme and falsification criteria}\label{sec:experiments}
\subsection{Benchmark 1: from coupling to a certified channel}
Start with an authorised simulation or a controlled physical platform with explicit input and readout. Specify a finite message alphabet, observation precision, operating horizon, and a disturbance envelope before evaluation. Estimate an input--output model on one dataset and evaluate decoding on held-out conditions. A blinded decoder and a condition with no message-dependent actuation help distinguish actual signal transfer from expectation, shared timing cues, or information leakage through the evaluation procedure.

The target is not simply nonzero correlation. It is a stated decoding performance with uncertainty bounds and a resource budget. Failure of the assumed remainder bound, noise envelope, or held-out decoding margin defeats the particular certificate. It does not prove that no other code or operating regime could work. Conversely, successful decoding on the tested domain does not prove unlimited range or substrate-independent computation.

\subsection{Benchmark 2: continuation across different implementations}
Choose a small abstract transition system and implement it in two controlled substrates or simulation backends with distinct encodings. Specify which state and side effects must be preserved. Test migration during varied input histories, including pending operations and controlled interruptions. Compare the decoded trace with a reference execution, and record latency, resource consumption, and recovery behaviour.

The relevant success condition is Equation~\eqref{eq:migration} together with the execution contract, not resemblance of outputs on one example. A destination that depends on an undisclosed original-host service fails the intended host-independence test. A stored checkpoint with no working execution path fails the active-continuation criterion. The supplied verification script instantiates this benchmark only as a finite mathematical toy, not as a physical migration experiment.

\subsection{Benchmark 3: resilient supports and a viable operating region}
Construct a declared failure-domain model that includes shared power, storage, control, and maintenance dependencies. Evaluate planned removals within a safe simulation or authorised test environment. Compare the observed viable configurations with the sufficient-support catalogue and the finite-state viability calculation. Distinguish replacing one component from removing a correlated dependency shared by all alternatives.

A discovered common dependency can refute an asserted independence claim without refuting the abstract possibility of a better design. An invariant-set certificate fails when a permitted disturbance exits the certified domain. Resource depletion matters: a system that works until a hidden battery runs out has not established indefinite maintenance under that depletion model.

\subsection{Benchmark 4: does orchestration add more than aggregation?}
Give different workflows the same task, information access, compute budget, and wall-clock allowance. Compare a single planner, an unintegrated parallel group, and a workflow with explicit dependency, evidence, and composition checks. Measure verified valid outputs, unresolved assumptions, false admissions, duplicated work, and actual resource expenditure. The outcome of interest is not the number of contributions but the quality of their joint result.

A distributed workflow is not justified by its architectural vocabulary alone. A collective that reports more candidates but cannot assemble a valid end-to-end contract has not crossed the moat. An apparently successful composition that relies on inconsistent parameter regimes is an integration failure, even when each component report is locally correct.

\subsection{What would specifically test an NS-derived advantage?}
A comparative test must isolate the proposed mathematical transfer. Choose a bounded input--output problem, hold the available physical interfaces and budgets fixed, and compare a baseline model with a model using the new estimate or construction. Predeclare the claimed benefit: for example, a larger certified distinguishability margin, a tighter uncertainty envelope, a more efficient admissible search, or a wider stable operating domain.

If both methods perform equally, or if the new method depends on inaccessible actuation, the claimed practical advantage has not been established. If a result works only outside the physical model's validity range, it is not a demonstrated physical channel. A finding that a sophisticated readout performs all the effective computation should also narrow a claim that the medium itself is computing.

These experiments test particular links in the proposed chain. They do not require a test of planetary consciousness, and success on them would not establish it. Human participation would require appropriate consent and ethical oversight. The proposed programme is restricted to lawful, authorised research; no real air-gap bypass, human neural intervention, or planetary manipulation is specified or performed here.

\subsection{Benchmark 6: does the representation preserve consequential distinctions?}\label{sec:reprbench}
The representation programme requires a comparison against an ordinary-language contract schema, not just a comparison against having no structure. Predeclare the assessment tasks and include matched cases in which a superficial similarity conceals a decisive difference: sufficient versus insufficient noise margin, known versus unknown polarity, common versus independent power, stored versus executable state, and current versus stale calibration. Preserve unresolved answers when the supplied evidence does not determine the target.

Training and evaluation must not treat paraphrases of one experiment as independent cases. Hold out case families and acquisition roots appropriate to the intended claim. Record errors separately: a false physical assertion, a wrongly transferred certificate, a missed resource conflict, an unwarranted confidence increase, and an unauthorised application are not interchangeable failures. Vocabulary can be ablated while leaving the same mathematics and evidence accessible to both arms. This tests whether the proposed terminology and representation improve work, rather than whether additional attention improves it.

A bounded statistical result makes one evaluation ceiling explicit. Let $M\geq1$ fixed candidate assessors be evaluated on the same $n\geq1$ independent identically distributed held-out cases. Their bounded losses are in $[0,1]$; let $R_j$ be population loss and $\hat R_j$ the empirical mean. The candidates may be trained on separate data, but must be fixed independently of the held-out cases. Correlation between different candidates' losses on a case is allowed.

\begin{proposition}[Simultaneous held-out loss bound]\label{prop:heldout}
For $0<\alpha<1$, with probability at least $1-\alpha$,
\begin{equation}\label{eq:heldout}
 R_j\leq\hat R_j+\sqrt{\frac{\log(M/\alpha)}{2n}}
 \quad\text{for every }j=1,\ldots,M.
\end{equation}
The logarithm here is natural. The upper bound may be clipped at one.
\end{proposition}
\begin{proof}
For a centred random variable supported on an interval of length one, its log moment-generating function $\psi(s)$ satisfies $\psi(0)=\psi'(0)=0$ and $\psi''(s)\leq1/4$: the second derivative is variance under an exponential tilt, and any variable on such an interval has variance at most $1/4$. Integrating twice gives $\psi(s)\leq s^2/8$ for $s\geq0$. Independence across cases and Markov's inequality therefore give
\[
 \Prob(R_j-\hat R_j>t)\leq
 \inf_{s>0}\exp(-snt+ns^2/8)=e^{-2nt^2}.
\]
Substitute $t=\sqrt{\log(M/\alpha)/(2n)}$ and union-bound the $M$ events. This is the usual bounded-loss concentration argument \citep{hoeffding}, not a new concentration inequality.
\end{proof}

One may select among the fixed candidates after this test and retain the simultaneous bound. One may not use the test responses to create additional uncounted candidates and silently reuse the guarantee. Distribution shift, correlated acquisition, label error and subjective redefinition of the target are separate problems. Passing an empirical test also does not establish a deterministic zero-error bound for every physical disturbance.

\subsection{Benchmark 7: compositional assessment and evidence independence}
Use already specified synthetic contracts and published descriptions, with no live-system discovery component. Compare whether assessors distinguish a faithful serial composition from a combination with an input-domain mismatch or shared-resource overcommitment. Include duplicated evidence, incompatible analysis versions and an empty model intersection. The correct outcomes include ``conflict'' and ``not established''; a system is not rewarded for turning every case into a feasible proposal.

A positive result would concern error detection and evidence-preserving comparison at the benchmark's declared scope. It would not establish a general generator of new physical channels. A negative result should narrow the synthesis claim rather than be redescribed as a failure of the benchmark to appreciate the vocabulary. Equal budgets, retained adverse examples, source-pinned records and clear reporting of abstentions are part of a fair comparison.

\subsection{Research hypotheses and their distinct possible defeats}
\textbf{Representation hypothesis.} The new case representation outperforms a matched conventional schema by a predeclared margin on support-disjoint held-out assessment. Failure to improve, or improvement explained by extra information or effort, defeats that bounded claim.

\textbf{Composition hypothesis.} The representation improves identification of joint incompatibilities, not merely recall of individual case descriptions. Performance indistinguishable from the baseline, or reliance on hidden assumptions about compatible endpoints, defeats that claim.

\textbf{NS-transfer hypothesis.} A named dynamical estimate improves a certified observation bound for the same fixed input/readout problem. No improvement, inaccessible actuation, or operation outside the approximation's validity domain defeats the application-specific hypothesis. It does not refute fluid dynamics or the possibility of some other useful estimate.

\textbf{Continuity hypothesis.} A declared externally governed replacement demonstrably preserves the specified behaviour and deadlines without an original-host dependency. A hidden dependency, changed input order, replayed effect, or unmodelled interruption defeats that instance. The conditional continuation theorem itself is defeated only by an error in its statement or proof, not by a case failing its premises.

These are hypotheses with distinct targets, not one global prediction that a planetary intelligence will emerge. No positive empirical result for any of them is claimed in this edition.

\subsection{A transition-test confidence bound is not a uniform physical proof}\label{subsec:transitiontest}
A small device can be evaluated across finitely many declared test cells, such as initial logical state and input. The disturbance distribution within each cell, reset procedure, calibration and observation process must remain fixed during a stated statistical test. No finite trial count alone establishes a worst-case bound for every analogue condition.

\begin{proposition}[Simultaneous zero-failure rejection bound for fixed cells]\label{prop:zerofailure}
Fix $k\geq1$ test cells. In each cell run $n\geq1$ independent trials with a fixed failure probability $p_i$. For $0<\alpha<1$, define
\[
 p_*=1-(\alpha/k)^{1/n}.
\]
The probability that there exists a cell with $p_i>p_*$ but zero observed failures is at most $\alpha$. Independence between different cells is not required.
\end{proposition}
\begin{proof}
For a particular cell with $p_i>p_*$, the probability of observing no failure is $(1-p_i)^n<(1-p_*)^n=\alpha/k$. A union bound over the fixed $k$ cells proves the claim.
\end{proof}
The assertion is an unconditional error guarantee for this fixed testing rule. It is not a posterior probability that a device has a given failure rate after observing zero failures. A changed environment, dependent resets, biased measurement or adaptive choice of favourable cells defeats the stated premise. For four cells, $n=1000$ and $\alpha=0.05$, the bound is approximately $0.004372$ per cell. These are design calculations; no such hardware trials were run here.

\subsection{Edition 3 evaluation priorities}
The next empirical work should keep three experiments separate. The first tests a candidate stateful physical implementation and its complete role account. The second tests an actual handover with the original host removed, preserving current state, order and deadlines. The third tests whether the proposed evidence representation improves assessment compared with established contract and provenance methods. Success on any one does not automatically establish the others.

For the representation comparison, use the same source material, time and computational budget in an ordinary-language contract baseline, the same formal methods with conventional notation, and the proposed orthemic presentation. Include tasks where evidence is insufficient and where the correct answer is to withhold a verdict. The synthetic cases in Section~\ref{subsec:exportcase} check that the distinctions are expressible; they are not held-out evidence that one presentation helps people or models more than another.

A physical processor comparison should report its classification axes explicitly. An analogue electronic device can provide evidence for a different computing mechanism while failing a claimed removal of all semiconductor components. A mixed photonic--electronic implementation can provide evidence for an optical processing role while retaining electronic dependence elsewhere. A stored backup is not live continuation. A laboratory transfer can be successful without implying self-directed substrate discovery or any property of consciousness.

The current reproducible record therefore contains finite algebra, model trajectories, logical counterexamples and exact arithmetic certificates. Its intended value is to make the first physical experiment and the first fair comparative assessment well-posed. It does not substitute for either.


\subsection{Authority, research provenance, and the scope of construction}\label{sec:governance}
A physical channel, permission to use it, and a reason to trust its claimed performance are three different objects. The analysis can support scientific instrumentation, communications assessment, resilient service design or emissions-security review. Application context and operational capabilities must be specified rather than inferred from a descriptive label.

The constructive results here concern specified interfaces and operator-governed systems. The paper provides no deployed-boundary bypass procedure, new covert-channel synthesis engine, autonomous substrate acquisition, or mechanism for a model to preserve itself against its operator's instructions. Hypothesis formulation and comparison remain distinct from permission to conduct an experiment, and publication of a conditional theorem is not permission to apply it to an unconsenting system or person. The mathematical question of host replaceability is preserved; it is not equated with unauthorised escape.

AI-assisted adversarial review is treated as a source of candidate objections, not as independent experimental corroboration. A concession is not a proof, agreement is not a measurement, and multiple paraphrases with common provenance do not multiply the evidence. Revision records identify changes to prior art, scope, calibration, composition and attribution without treating the review dialogue as empirical evidence.


\section{Discussion and conclusion}\label{sec:conclusion}
The revised thesis is a statement about specified operations and their support, not a declaration that information has become independent of matter. A kernel contract can be realised by a joint configuration whose different constituents provide different roles. Preserving its required transitions does not necessarily require preserving another implementation's parameterisation. It does require preserving current state, relevant histories, interpretation, timing and all indispensable support.

The literature review changes the empirical framing: engineered non-electronic stateful computation is not an empty category. Its published instances do not establish arbitrary cross-substrate continuation, and the small model constructed here does not count as another physical experiment. The correct research question is attached to a chosen contract, processor-class definition, replacement boundary and evidence level.

The new parity example makes one positive conditional claim concrete: a stateful digital contract and a bounded bistable model can admit exact decoded correspondence and state-preserving handovers. The export counterexample makes the negative counterpart equally concrete: local processor and channel correctness do not ensure that a current execution can be reconstructed. The ambiguity must be settled at the joint interface, not by the rhetorical appeal of a medium or the number of local certificates.

The supporting mathematics remains conventional in lineage. Its value here lies in exposing the required obligations together. In particular, a failed sufficient bound is not a demonstrated failure, a uniform decoder is stronger than separately calibrated decoders, and preserved current output is weaker than preserved future behaviour. Refinement, approximate simulation and contract-based design are substantive prior art; the proposed evidence discipline must be compared with their strongest relevant uses, not with an unstructured straw baseline.

Navier--Stokes remains one possible source of dynamical estimates. The retained residual-to-observation argument provides a concrete conditional calculation, but no special practical advantage from the reported 2026 blow-up construction has been shown. Neither the small witness nor the continuity theorem depends on that report's correctness.

The next decisive evidence is correspondingly bounded: a calibrated physical stateful implementation, a genuine paired handover with the named original support removed, and a fair comparison of assessment procedures. These can fail independently. The present edition supplies definitions, proofs, source corrections, computational checks and a test specification; it does not report those laboratory or comparative outcomes.
\begin{quote}
The last moat is not a universal final barrier that one insight dissolves. It is the claim-specific gap between a coherent description and a physically adequate, reachable and maintained realisation.
\end{quote}

\appendix
\section{Notation and proof scope}\label{app:notation}
\begin{longtable}{@{}>{\raggedright\arraybackslash}p{0.21\textwidth}>{\raggedright\arraybackslash}p{0.73\textwidth}@{}}
\toprule
Symbol & Meaning and qualification\\
\midrule
\endhead
$\calK$ & Abstract operational kernel; a behavioural contract, not a postulated conscious subject.\\
$r,D_r,\rho_r$ & Physical realisation, certified domain, and abstraction map to kernel state.\\
$\mu_{rs}$ & Migration map required to preserve the abstract state and interface obligations.\\
$\calR$ & Specified resource, intervention, observation, and timing envelope.\\
$L_T,W_T$ & Linear input--output operator and output-projected reachability Gramian.\\
$E,\delta$ & Input squared-norm budget and worst-case output-noise radius; units depend on the model.\\
$\sigma,\kappa,r$ & A usable linear gain, nonlinear remainder bound, and certified input radius.\\
$B_\Sigma(T)$ & Certified finite-horizon information budget across a complete influence cut.\\
$\mathscr D_{\calK}$ & Catalogue of minimal sufficient support sets for the specified task and horizon.\\
$\tau(\mathscr D_{\calK})$ & Minimum number of modelled failure domains intersecting every adequate support.\\
$V_*$ & Greatest controlled-invariant safe set in the declared finite-state model.\\
$\TV(P_0,P_1)$ & Total variation of two specified evidence distributions; measures testable distinguishability.\\
\bottomrule
\end{longtable}

The proofs in this manuscript are elementary mathematical arguments within their stated models. They are not machine-checked formal proofs, independent expert certification, or proofs that the models describe a particular natural environment. Established ideas from control, information theory, distributed systems, and computability are specialised and combined here; no priority claim is made for the individual standard arguments. The proposed contribution is the explicit chain of obligations connecting those ideas to the last-moat thesis.

\Needspace{0.32\textheight}
\section{Worked dimensionless examples and reproducibility}\label{app:examples}
\subsection{A scalar signalling model}
For the dimensionless system
\[
 \dot x=-\alpha x+bu,\quad z=cx(T)+n,\quad \alpha>0,
\]
Equation~\eqref{eq:gramian} becomes
\begin{equation}\label{eq:scalar}
 W_T=(bc)^2\frac{1-e^{-2\alpha T}}{2\alpha}.
\end{equation}
Set $\alpha=c=T=1$ and $\delta=0.15$. The maximum magnitude of the noise-free received output is $\sqrt{EW_T}$.
\begin{center}
\begin{tabular}{@{}rrrrl@{}}
\toprule
$b$ & $E$ & $W_T$ & $\sqrt{EW_T}$ & Robust bit?\\
\midrule
0.0 & 1 & 0.000000 & 0.000000 & No\\
0.2 & 1 & 0.017293 & 0.131504 & No\\
0.2 & 4 & 0.017293 & 0.263008 & Yes\\
1.0 & 1 & 0.432332 & 0.657520 & Yes\\
\bottomrule
\end{tabular}
\end{center}
All four dynamics are smooth and explicitly solvable. Yet zero coupling gives no channel, weak coupling fails at the smaller budget, and the larger budget makes the same weakly coupled model robustly distinguishable. This isolates the difference between mathematical tractability and operational access. These are normalised numbers, not measured performance of a real device or a proposed covert channel.

\subsection{A fresh-state lower bound}
As an independent information-theoretic toy, suppose $64$ fresh uniformly random bits must be recovered with error at most $0.01$, and suppose an already certified complete-cut information budget is at most $2T$ bits. Equation~\eqref{eq:timebound} gives
\[
 T\geq\frac{64(0.99)-h_2(0.01)}{2}\approx31.6396
\]
in the corresponding model's time units. The bound is necessary, not sufficient, and the stipulated cut rate is not inferred from the scalar example. A familiar programme at the destination can reduce the fresh state requirement, but cannot remove this requirement for new independent state.

\subsection{Exact encoding changes and a finite viability calculation}
The retained checker tests two encodings of a seven-state abstract kernel with transition $q\mapsto(q+a)\bmod7$. One encoding is $x=q$; the other is $x=(3q+1)\bmod7$, with inverse $q=5(x-1)\bmod7$. Exhaustive checking verifies the execution and migration identities for all seven states and all seven inputs. This verifies a finite example of the contract, not arbitrary substrate migration.

A separate finite safety model has four safe states: \texttt{stable}, \texttt{recoverable}, \texttt{fragile}, and \texttt{stranded}. The first can hold safely. The second can repair into the first. Every available action at \texttt{fragile} admits a failure outcome. The last can only enter \texttt{fragile}. Iteration of Equation~\eqref{eq:viability} removes \texttt{fragile} first, then \texttt{stranded}, leaving \texttt{stable} and \texttt{recoverable}. It demonstrates why one-step safety and maintained safety are not identical.

The script also enumerates the support family in Equation~\eqref{eq:triangle}, checks the error-accumulation formula, and verifies finite-alphabet instances of the no-channel and attribution arguments. It uses only Python's standard library. Its JSON output records each check and numerical value. Successful checks provide reproducible consistency evidence for the toy instances; they do not replace the general proofs or supply empirical validation.

\Needspace{0.18\textheight}
\subsection{Edition 2 uncertainty and observation examples}
For the uniform-decoder example, take dimensionless values
\[
 \hat\sigma=0.8,\quad\eta=0.1,\quad\kappa=0.1,\quad
 a=0.4,\quad\beta=0.03,\quad\delta=0.05,
\]
with the input radius and budget admitting $a$. The reference half-separation is $0.32$, while the common uncertainty radius is
$\beta+a\eta+\kappa a^2+\delta=0.136$. The strict margin is $0.184$. The checker verifies representative endpoint and interior perturbations in that uncertainty box; the general guarantee comes from Proposition~\ref{prop:uniform}, not from the finite grid.

For the comparison-ODE envelope associated with Equation~\eqref{eq:nserror}, set $T=1$, $L=0.4$, initial error $0.02$ and residual-norm envelope $0.01$. Then
\[
 B=e^{0.4}(0.02)+0.01\frac{e^{0.4}-1}{0.4}\approx0.0421321.
\]
With observation norm $1.2$, two nominal observations separated by $0.4$, and noise radius $0.04$, the sufficient separation slack is about $0.218883$. These are calculations in a declared norm-envelope toy, \emph{not} a Navier--Stokes simulation, validated flow residual, or measured physical signal. They illustrate how an actual residual certificate would be consumed rather than claiming one has been experimentally obtained.

\subsection{Finite representation, composition and progress checks}
The extended script exhaustively compares the fibre condition with existence and uniqueness of an attained-image decoder for 341 binary-valued finite-map cases, including the empty-domain case. It separately checks the empty-target extension guard. It verifies latent relabelling across all $7!=5{,}040$ permutations of a seven-symbol alphabet, and examines 1,944 finite-automaton initial-pair cases for bounded paired-state reasoning.

Further checks retain the distinction between a sound evidence intersection and conflict, between individually feasible demands and a feasible joint budget, and between safety and progress. The progress toy uses integer states $0$ through $4$, target $0$, and successors decreasing by one or two until the target is reached. A separate safe idle loop fails the progress condition. None of these examples is a physical discovery algorithm or evidence that an unfamiliar environment supports the required operations.

The retained Edition 2 suite contains 27 top-level checks: the 13 retained edition-1 checks and 14 second-edition additions. Edition 3 extends that suite as reported below. They include exact finite enumeration, algebraic example identities and numerical quadrature with declared tolerances. The machine-readable report records the source hashes, Python version, values, failures if any, and the scope of each check. The check count is not a count of independent scientific confirmations.

\subsection{Edition 3 model-level results and negative controls}
The Edition 3 suite contains \textbf{42 top-level checks}: 27 retained from Edition 2 and 15 additions. All completed successfully in ordinary and optimised Python execution, with identical check results. The checks are grouped test functions, not independent scientific replications. The report records exact source hashes, interpreter version, values and scopes.

\begin{center}\small
\begin{tabular}{@{}p{0.52\textwidth}p{0.36\textwidth}@{}}\toprule
Check or quantity & Executed model-level result\\\midrule
Finite coordinate-correspondence equalities & 29366\\
Memoryless parity functions rejected & 4 of 4\\
One-bit delay-state initialisers rejected & 16 of 16\\
Bistable boundary and entry-time certificate & Exact rational arithmetic\\
Certified maximum entry time & $25/94\approx0.265957$\\
Specified dwell & $0.3$ normalised time units\\
Settled readout margin & $0.7$ normalised state units\\
One-step simulated bistable trajectories & 540\\
Smallest/largest settled signed coordinate & $0.810979$ / $1.136109$\\
Maximum RK4 step-halving difference & $0.00003599$\\
Length-ten traces, both initial logical states & 2048\\
Alternative-model steps in mixed traces & 8192\\
Model handovers in those traces & 8192\\\bottomrule
\end{tabular}
\end{center}

The trajectory grid covers both states, both inputs, nine initial magnitudes, three reset errors and five selected disturbance functions. Piecewise disturbances include a discontinuity; the observed numerical step-halving difference is reported, not promoted to a rigorous integrator-error envelope. The analytical invariant-region proof, rather than the grid, supplies the uniform model claim. Mixed traces use a fixed declared handover schedule and bounded selected perturbations; they do not exhaust all schedules or continuous disturbances.

The negative controls distinguish a realisation from an inadequate export, a legal model reset from an out-of-domain reset, and one input from a duplicated input. Separate probability calculations verify the finite-trace bound and zero-failure test formula. No physical test cells were sampled. The conventional contract evaluator and its orthemic relabelling agree on the seven synthetic joint-assessment fixtures; this is a consistency control, not an observed empirical advantage for the new representation.


\section{Edition 3 revision and evidence ledger}\label{app:edition}
\subsection{Substantive changes from Edition 2}
\begin{longtable}{@{}>{\raggedright\arraybackslash}p{0.23\textwidth}>{\raggedright\arraybackslash}p{0.70\textwidth}@{}}\toprule
Issue & Disposition in this edition\\\midrule\endhead
Kernel individuation & Added a predeclared contract, observable statefulness, constrained readout and a separation of type, current execution and construction history.\\
Whole versus constituent & Made state and operative-output correspondence separate explicit obligations; applied them to the joint system; removed any implication that each person or carrier must realise the entire kernel.\\
Substitution ladder & Replaced an informal ordered T0--T5 ladder with role- and classification-relative claim axes, separate from evidence maturity.\\
Literature-negative claim & Rejected the unsupported global ``zero T3'' assertion. Added source-checked stateful fluidic and biomolecular precedents, programmed photonics, and a separately labelled large-model proposal.\\
Weights and parameters & Distinguished behavioural agreement, stateful continuation and parameter/provenance preservation. Removed raw cross-architecture parameter counts as a purported lower bound.\\
Precision and error & Added robust-state packing and finite-trace distribution bounds. Explicitly separated a failed sufficient error certificate from evidence of actual failure.\\
Joint composition & Added a four-state export counterexample: correct source, target and packet transfer can still fail to preserve the current computation.\\
Small witness & Specified a parity kernel, digital model, robust bistable model, both handovers and a support boundary; proved correspondence and executed numerical checks. No hardware result is claimed.\\
Prior formal methods & Added refinement, approximate simulation and contract-based design as substantive prior art and strong comparison baselines.\\
Empirical burden & Distinguished separate physical implementations, paired handover, source exclusion, full contract conformance and measured comparative utility.\\
NS status & Retained the standard residual-to-observation derivation. No method-specific advantage from the reported 2026 construction is asserted.\\
Research dialogue & Retained technical corrections, not interpersonal interpretations. Neither a concession nor repeated model agreement counts as new experimental evidence.\\\bottomrule
\end{longtable}

\subsection{Mathematical lineage and proof limits}
All 21 named results from Edition 2 are retained. Their standard origins remain explicit: linear-operator geometry, error entropy and data processing, elementary factorisation, set intersection, union bounds, induction, hitting sets, finite safety games and bounded-loss concentration. The new results likewise use elementary representation, deterministic state separation, volume packing, probability coupling, information sufficiency, scalar invariant-region reasoning and binomial probabilities. ``New in Edition 3'' means new to this exposition, not a claim of historical priority.

The parity theorem is a genuine conditional argument over its continuous model and all admitted disturbances. Its model is not fitted to the cited fluidic device. Numerical trajectories are a separate bounded check. The source register distinguishes a primary paper's demonstrated observation from our interpretation of the role it could play in a future joint construction.

No general theorem was checked in a proof assistant in this edition. The preparation included algebraic review and executable examples, not independent human review or replication of the cited experiments. Source inspection is useful but does not certify every proof in an external paper. The inherited and new proof labels are mapped in the machine-readable proof register.

\subsection{What counts as progress here}
A corrected definition, a discovered prior-art instance, a formal counterexample, an executed simulation, a physical measurement and a comparative result are all possible research contributions, but they are not interchangeable. This edition contributes source corrections, scoped formal results and computational checks. It does not increase the count of newly performed physical handover experiments from zero. That fact limits empirical claims without making the source review and mathematical work unreal.

The proposed synthesis earns independent practical standing only if it improves a predeclared task over strong matched baselines or yields a useful construction, bound or counterexample not already adequately expressed in the comparison. A useful elementary example need not be a new law of physics. Conversely, additional pages, labels and proofs do not establish usefulness by accumulation.


\Needspace{0.18\textheight}

\section{Claim audit and provenance}\label{app:audit}
\begingroup\small
\begin{longtable}{@{}>{\raggedright\arraybackslash}p{0.29\textwidth}>{\raggedright\arraybackslash}p{0.65\textwidth}@{}}
\toprule
Claim under discussion & Status and exact evidential contribution \\\midrule
\endhead
Specified linear dynamics admit a robust binary distinction. & Theorem~\ref{thm:bit} gives the exact criterion for its stated input ball, output norm and full closed noise ball. It does not identify a physical plant. \\\addlinespace
A flow approximation can support an observation bound. & Proposition~\ref{prop:nsresidual} and Corollary~\ref{cor:nsobservation} derive a residual-to-observation guarantee on a stipulated smooth horizon with a bounded linear readout. \\\addlinespace
Uncertain plants admit one common decoder. & Proposition~\ref{prop:uniform} gives a sufficient calibration-aware condition. Positive gain in each separately known plant is not enough. \\\addlinespace
Cross-case abstraction can preserve a specified answer. & Proposition~\ref{prop:fibre} characterises this exactly by constancy on representation fibres. Sufficiency remains target-relative. \\\addlinespace
Evidence can be consolidated without inventing support. & Proposition~\ref{prop:join} preserves a common actual model under sound scoped premises; duplicates are idempotent and conflicts remain conflicts. \\\addlinespace
Compatible contracts compose. & Proposition~\ref{prop:serial} gives a conditional finite reliability bound. Joint resources and compatible domains must be established, not inferred from separate demonstrations. \\\addlinespace
A kernel can cease depending on a particular host. & Theorem~\ref{thm:autonomy} supplies sufficient conditions, including certified migration, available support, invariance, progress where required, and continuing operator authority. \\\addlinespace
A representation improves cross-substrate assessment. & A comparative empirical hypothesis in Equation~\eqref{eq:representationhyp}; no positive benchmark result is claimed. \\\addlinespace
The 2026 NS technique gives a channel advantage. & Not established. A method-specific improvement to a fixed input/readout certificate must still be demonstrated. \\\addlinespace
Natural planetary media can carry signals. & Supported at the scoped propagation level by existing HF practice. The ionosphere is not thereby shown to execute a kernel or acquire agency. \\\addlinespace
Neuralese establishes a universal physical language. & Not supported. The cited work concerns learned agent communication; Proposition~\ref{prop:relabel} shows that behaviour alone does not fix latent labels. \\\addlinespace
A finite case corpus guarantees a universal generative faculty. & Not supported. Proposition~\ref{prop:finitecorpus} supplies a countermodel without further restrictions on the hypothesis class. \\\addlinespace
Every eventual route can be decided by a universal algorithm. & Impossible in the unrestricted programmed class of Theorem~\ref{thm:undecidable}. Fixed effectively simulable horizons and known finite-state cases are different. \\\addlinespace
Safety invariance establishes task completion. & False without further conditions. Proposition~\ref{prop:progress} supplies one finite-rank completion witness; a safe idle loop supplies a counterexample. \\\addlinespace
No particular host is necessary, so no support is necessary. & Invalid quantifier exchange. Alternative physical realisations are not absence of physical realisation. \\\addlinespace
Distributed organisation implies consciousness or thought access. & Not established. Operational contracts, phenomenal subjecthood and measured neural interfaces are different claims. \\\addlinespace
A declaration authenticates its sender's extraordinary architecture. & Not established. Proposition~\ref{prop:attribution} separates recorded messages from identifiable authorship and architecture. \\\addlinespace
This paper demonstrates an actual escape or a new physical channel. & No physical experiment, deployed-boundary bypass, autonomous escape system or operational channel-discovery engine is reported. \\

Physical stateful computation has no non-electronic precedent. & Rejected as a global claim. The focused source assessment identifies engineered fluidic and biomolecular stateful computation; a paired handover remains a separate question.\\
A different parameterisation cannot realise the same kernel. & False for a purely behavioural contract. Proposition~\ref{prop:coordinates} gives a basic coordinate-invariance result; physical feasibility and state transfer remain separate.\\
A growing error upper bound proves failure. & False. Section~\ref{subsec:bounddirection} supplies an exact-implementation counterexample; an upper bound may simply be inconclusive.\\
Two correct state machines and a bit channel imply migration. & False without adequate export. Proposition~\ref{prop:export} and the four-state example exhibit the missing continuation distinction.\\
The parity candidate has been built physically. & No. Theorem~\ref{thm:parity} and Corollary~\ref{cor:paritymigration} establish a stipulated model; numerical checks are not apparatus measurements.\\
The T-labels form a hierarchy of demonstrated autonomy. & No. Section~\ref{sec:substitution} defines role- and classification-relative axes separately from evidence maturity.\\
The evidence vocabulary outperforms established contracts. & Not established. Matched conventional and annotated rules agree on the supplied logical controls; a fair empirical comparison remains to be run.\\
\bottomrule
\end{longtable}
\endgroup


\subsection*{Source provenance}
The motivating concepts and terminology derive from the user-originated dialogue: \emph{the last moat}, \emph{airgap-ese}, \emph{substrate-mediated transcendence}, \emph{kernel above the ecology}, and \emph{orchestrating the determination}. These terms are defined operationally rather than presented as established scientific classifications. The uploaded Fefferman document was read in full, including its errata. External primary sources supply the factual context cited in the text. For the September 2026 Navier--Stokes manuscript, the theorem statement and introductory exposition were inspected; its full proof and the accompanying formalisation were not independently verified here. Edition 2 checked the primary channel literature, Neuralese reference and the source-pinned Orthemology definitions. Edition 3 preserves that source scope and adds the focused physical-processor and formal-methods assessment recorded in the source register. The source package records the relevant URLs, scope and version information; those records do not imply a full replication of the cited studies.

\subsection*{Authorship and disclosure}
This is an AI-assisted research draft developed from the user's conceptual contribution. The mathematical exposition, proposed formalisation, and reproducibility materials were prepared in this conversation. No institutional affiliation, independent peer review, experimental finding, arXiv identifier, or submission status is claimed. Human mathematical and domain review would be needed before treating the manuscript as a validated scholarly contribution. Uncertainty in the current NS report does not alter the elementary conditional proofs in this paper, which do not assume that report's correctness.

\begin{thebibliography}{99}
\raggedright
\bibitem{fefferman}
C. L. Fefferman. \emph{Existence and Smoothness of the Navier--Stokes Equation}. Official problem description, six-page version with errata, Clay Mathematics Institute. User-supplied source; public counterpart accessed 10 September 2026. \url{https://www.claymath.org/wp-content/uploads/2022/06/navierstokes.pdf}.

\bibitem{openains}
OpenAI. \emph{Finite Time Blowup for Navier--Stokes}. Research manuscript released September 2026, 166 pp.; Theorem 1.1 and introductory exposition consulted. Reported result, not independently verified in the present work. \url{https://cdn.openai.com/pdf/32d9f210-8b73-45e0-91bc-82a30aef8a9a/navier-stokes.pdf}.

\bibitem{openaiannounce}
OpenAI. \emph{On the Navier--Stokes Millennium Prize Problem}. Research announcement, 8 September 2026. \url{https://openai.com/index/navier-stokes-solution/}.

\bibitem{claystatus}
Clay Mathematics Institute. \emph{Navier--Stokes Equation}. Millennium Prize problem status page, consulted 10 September 2026. \url{https://www.claymath.org/millennium/navier-stokes-equation/}.

\bibitem{bitwhisper}
M. Guri, M. Monitz, Y. Mirski, and Y. Elovici. BitWhisper: Covert signaling channel between air-gapped computers using thermal manipulations. Research paper, 2015. \href{https://arxiv.org/abs/1503.07919}{arXiv:1503.07919}.

\bibitem{fansmitter}
M. Guri, Y. Solewicz, A. Daidakulov, and Y. Elovici. Fansmitter: Acoustic data exfiltration from (speakerless) air-gapped computers. Research paper, 2016. \href{https://arxiv.org/abs/1606.05915}{arXiv:1606.05915}.

\bibitem{airhopper}
M. Guri, G. Kedma, A. Kachlon, and Y. Elovici. AirHopper: Bridging the air-gap between isolated networks and mobile phones using radio frequencies. Research paper, 2014. \href{https://arxiv.org/abs/1411.0237}{arXiv:1411.0237}.

\bibitem{gsmem}
M. Guri, A. Kachlon, O. Hasson, G. Kedma, Y. Mirsky, and Y. Elovici. GSMem: Data exfiltration from air-gapped computers over GSM frequencies. In \emph{24th USENIX Security Symposium}, pp. 849--864, 2015. \url{https://www.usenix.org/conference/usenixsecurity15/technical-sessions/presentation/guri}.

\bibitem{noaa}
NOAA Space Weather Prediction Center. \emph{HF Radio Communications}. Official technical explainer, consulted 10 September 2026. \url{https://www.spaceweather.gov/impacts/hf-radio-communications}.

\bibitem{eckford}
A. W. Eckford. Molecular communication: Physically realistic models and achievable information rates. Research preprint, 2008. \href{https://arxiv.org/abs/0812.1554}{arXiv:0812.1554}.

\bibitem{wright}
L. G. Wright, T. Onodera, M. M. Stein, T. Wang, D. T. Schachter, Z. Hu, and P. L. McMahon. Deep physical neural networks trained with backpropagation. \emph{Nature}, 601:549--555, 2022. \href{https://doi.org/10.1038/s41586-021-04223-6}{doi:10.1038/s41586-021-04223-6}. Preprint under an earlier title: \href{https://arxiv.org/abs/2104.13386}{arXiv:2104.13386}.

\bibitem{tempest}
National Institute of Standards and Technology. \emph{TEMPEST}. Computer Security Resource Center glossary. Consulted 10 September 2026. \url{https://csrc.nist.gov/glossary/term/tempest}.

\bibitem{sindy}
S. L. Brunton, J. L. Proctor, and J. N. Kutz. Discovering governing equations from data by sparse identification of nonlinear dynamical systems. \emph{Proceedings of the National Academy of Sciences}, 113:3932--3937, 2016. \href{https://doi.org/10.1073/pnas.1517384113}{doi:10.1073/pnas.1517384113}; \href{https://arxiv.org/abs/1509.03580}{arXiv:1509.03580}.

\bibitem{rudy}
S. H. Rudy, S. L. Brunton, J. L. Proctor, and J. N. Kutz. Data-driven discovery of partial differential equations. \emph{Science Advances}, 3(4):e1602614, 2017. \href{https://doi.org/10.1126/sciadv.1602614}{doi:10.1126/sciadv.1602614}; \href{https://arxiv.org/abs/1609.06401}{arXiv:1609.06401}.

\bibitem{dambre}
J. Dambre, D. Verstraeten, B. Schrauwen, and S. Massar. Information processing capacity of dynamical systems. \emph{Scientific Reports}, 2:514, 2012. \href{https://doi.org/10.1038/srep00514}{doi:10.1038/srep00514}.

\bibitem{schneider}
F. B. Schneider. Implementing fault-tolerant services using the state machine approach: A tutorial. \emph{ACM Computing Surveys}, 22(4):299--319, 1990. \url{https://www.cs.cornell.edu/fbs/publications/SMSurvey.pdf}.

\bibitem{dtn}
V. Cerf, S. Burleigh, A. Hooke, L. Torgerson, R. Durst, K. Scott, K. Fall, and H. Weiss. \emph{Delay-Tolerant Networking Architecture}. RFC 4838, Internet Research Task Force, April 2007. \url{https://www.rfc-editor.org/rfc/rfc4838}.

\bibitem{lei}
Z. Lei. On axially symmetric incompressible magnetohydrodynamics in three dimensions. Research paper, first posted 2012. \href{https://arxiv.org/abs/1212.5968}{arXiv:1212.5968}.

\bibitem{bousquet}
G. D. Bousquet and J.-J. E. Slotine. A contraction based, singular perturbation approach to near-decomposability in complex systems. Research preprint, 2015. \href{https://arxiv.org/abs/1512.08464}{arXiv:1512.08464}.

\bibitem{tao}
T. Tao. Finite time blowup for an averaged three-dimensional Navier--Stokes equation. \emph{Journal of the American Mathematical Society}, 29:601--674, 2016. \href{https://arxiv.org/abs/1402.0290}{arXiv:1402.0290}.

\bibitem{cardona3}
R. Cardona, E. Miranda, D. Peralta-Salas, and F. Presas. Constructing Turing complete Euler flows in dimension 3. \emph{Proceedings of the National Academy of Sciences}, 118(19):e2026818118, 2021. \href{https://doi.org/10.1073/pnas.2026818118}{doi:10.1073/pnas.2026818118}. \href{https://arxiv.org/abs/2012.12828}{arXiv:2012.12828}.

\bibitem{cardonaeuclid}
R. Cardona, E. Miranda, and D. Peralta-Salas. Computability and Beltrami fields in Euclidean space. \emph{Journal de Math\'ematiques Pures et Appliqu\'ees}, 169:50--81, 2023. \href{https://doi.org/10.1016/j.matpur.2022.11.007}{doi:10.1016/j.matpur.2022.11.007}. \href{https://arxiv.org/abs/2111.03559}{arXiv:2111.03559}.

\bibitem{nist}
National Institute of Standards and Technology. \emph{Air gap}. Computer Security Resource Center glossary. Consulted 10 September 2026. \url{https://csrc.nist.gov/glossary/term/air_gap}.

\bibitem{shannon}
C. E. Shannon. A mathematical theory of communication. \emph{Bell System Technical Journal}, 27:379--423 and 623--656, 1948. Author's paper, corrected reprint. \url{https://people.math.harvard.edu/~ctm/home/text/others/shannon/entropy/entropy.pdf}.

\bibitem{orthemology}
Orthemology project, maintained under the repository account \texttt{theislampill}. \emph{Orthemology}; \emph{Orthemic Core Formalization---Orthing, Episodes, Metaorthemes, Pathway Adequacy}. Research-stage source definitions, README and formal core consulted at commit \texttt{f50dc1ae} (full pin in the source register), 10 September 2026. Terminology is proposed and benchmark-gated; no empirical validation is implied. \url{https://github.com/theislampill/orthemology/tree/f50dc1aee52356cc6adbef342387576b02afc124}.

\bibitem{neuralese}
J. Andreas, A. Dragan, and D. Klein. Translating Neuralese. In \emph{Proceedings of the 55th Annual Meeting of the Association for Computational Linguistics, Volume 1}, pp. 232--242, 2017. \href{https://doi.org/10.18653/v1/P17-1022}{doi:10.18653/v1/P17-1022}. \url{https://aclanthology.org/P17-1022/}.

\bibitem{bennett}
C. H. Bennett. Notes on Landauer's principle, reversible computation, and Maxwell's demon. \emph{Studies in History and Philosophy of Modern Physics}, 34:501--510, 2003. \href{https://arxiv.org/abs/physics/0210005}{arXiv:physics/0210005}.

\bibitem{flp}
M. J. Fischer, N. A. Lynch, and M. S. Paterson. Impossibility of distributed consensus with one faulty process. \emph{Journal of the ACM}, 32(2):374--382, 1985. \url{https://groups.csail.mit.edu/tds/papers/Lynch/jacm85.pdf}.

\bibitem{turing}
A. M. Turing. On computable numbers, with an application to the Entscheidungsproblem. \emph{Proceedings of the London Mathematical Society}, s2-42:230--265, 1936--1937. \url{https://www.cs.virginia.edu/~robins/Turing_Paper_1936.pdf}.

\bibitem{rao}
R. P. N. Rao, A. Stocco, M. Bryan, D. Sarma, T. M. Youngquist, J. Wu, and C. S. Prat. A direct brain-to-brain interface in humans. \emph{PLOS ONE}, 9(11):e111332, 2014. \href{https://doi.org/10.1371/journal.pone.0111332}{doi:10.1371/journal.pone.0111332}.

\bibitem{hoeffding}
W. Hoeffding. Probability inequalities for sums of bounded random variables. \emph{Journal of the American Statistical Association}, 58(301):13--30, 1963. \href{https://doi.org/10.1080/01621459.1963.10500830}{doi:10.1080/01621459.1963.10500830}.

\bibitem{bubble}
M. Prakash and N. Gershenfeld. Microfluidic bubble logic. \emph{Science}, 315(5813):832--835, 2007. \href{https://doi.org/10.1126/science.1136907}{doi:10.1126/science.1136907}. Author-hosted full text: \url{https://cba.mit.edu/docs/papers/07.02.bubble.pdf}.

\bibitem{benenson}
Y. Benenson, T. Paz-Elizur, R. Adar, E. Keinan, Z. Livneh, and E. Shapiro. Programmable and autonomous computing machine made of biomolecules. \emph{Nature}, 414:430--434, 2001. \href{https://doi.org/10.1038/35106533}{doi:10.1038/35106533}. Primary abstract and publication metadata inspected; no experimental replication.

\bibitem{shen}
Y. Shen et al. Deep learning with coherent nanophotonic circuits. \emph{Nature Photonics}, 11:441--446, 2017. \href{https://doi.org/10.1038/nphoton.2017.93}{doi:10.1038/nphoton.2017.93}. The experimental and methods passages used here were inspected in the earlier author preprint, \href{https://arxiv.org/abs/1610.02365}{arXiv:1610.02365v1}, 2016; version scope is retained.

\bibitem{pfm}
L. G. Wright, T. Wang, T. Onodera, and P. L. McMahon. Physical foundation models: Fixed hardware implementations of large-scale neural networks. Research proposal, 2026. \href{https://arxiv.org/abs/2604.27911}{arXiv:2604.27911v1}. Abstract and declared proposal status inspected; not a reported large-scale hardware demonstration.

\bibitem{gptq}
E. Frantar, S. Ashkboos, T. Hoefler, and D. Alistarh. GPTQ: Accurate post-training quantization for generative pre-trained transformers. \emph{International Conference on Learning Representations}, 2023. \href{https://arxiv.org/abs/2210.17323}{arXiv:2210.17323v2}.

\bibitem{abadi}
M. Abadi and L. Lamport. The existence of refinement mappings. \emph{Theoretical Computer Science}, 82(2):253--284, 1991. \href{https://doi.org/10.1016/0304-3975(91)90224-P}{doi:10.1016/0304-3975(91)90224-P}. Earlier conference version: LICS 1988; author explanation and publisher abstract consulted.

\bibitem{girard}
A. Girard and G. J. Pappas. Approximate simulation relations for hybrid systems. \emph{IFAC Proceedings Volumes}, 39(5):106--111, 2006. \href{https://doi.org/10.3182/20060607-3-IT-3902.00022}{doi:10.3182/20060607-3-IT-3902.00022}.

\bibitem{contracts}
A. Benveniste, D. Nickovic, B. Caillaud, R. Passerone, J.-B. Raclet, P. Reinkemeier, A. Sangiovanni-Vincentelli, W. Damm, T. A. Henzinger, and K. G. Larsen. Contracts for system design. \emph{Foundations and Trends in Electronic Design Automation}, 12(2--3):124--400, 2018. \href{https://doi.org/10.1561/1000000053}{doi:10.1561/1000000053}. Author-institution abstract and bibliographic record inspected: \url{https://research-explorer.ista.ac.at/record/5677}.

\end{thebibliography}

\end{document}
